% !TeX root = ../../Book4.tex
% 纲要标题：## 第11章　可积系统与作用角坐标
% 纲要位置：几何学计划纲要.md，第 3881 行。
% 正文状态：初稿；尚未正式审核。

% 纲要内容开始
% 
% > 本章目标与证明范围：将 Liouville–Arnold 的局部环面理论与矩映射凸性分类分章。完整证明紧连通正则共同水平集的环面结构和邻域坐标，随后展示整体单值障碍。
% 
% 纲要内容结束

\chapter{可积系统与作用角坐标}
\label{g4:ch:11}
\BookChapterNumber{g4:ch:11}

\begin{chapterabstract}
足够多的独立对合积分把 Hamilton 运动限制在半维共同水平集上。若其中一个正则连通分支紧，则交换流使它成为 Lagrange 环面；进一步把周期基延拓到邻近纤维，并积分典范一形式，便得到作用角坐标。本章完整证明这一局部结构，计算角变量的线性运动及坐标的整数仿射变换。随后沿底空间延拓周期格，用显式焦点型粘合模型展示非平凡单值怎样阻止全局坐标。振子、旋转曲面的测地流与球摆给出从正则环面到奇异纤维的具体计算。
\end{chapterabstract}

\begin{learninggoals}
\begin{itemize}
	\item 从微分独立性和 Poisson 对合证明正则共同水平集是 Lagrange 子流形。
	\item 由完备交换流构造实向量群作用，计算稳定子，并证明紧连通正则纤维的环面结构。
	\item 完整建立紧纤维的作用角邻域，说明周期基、作用积分及角变量各自的构造作用。
	\item 核对角周期为 \(2\pi\) 时的辛形式、频率与坐标变换，并区分闭轨道和稠密轨道。
	\item 建立周期格局部系统及整仿射结构，计算一个焦点型模型的单值矩阵。
	\item 计算振子、测地流和球摆的交换积分、正则纤维与临界行为，并区分二次近似和实际奇纤维。
\end{itemize}
\end{learninggoals}

两个互不耦合的振子各有一条能量椭圆。固定两个能量后，相空间中的共同水平集是两条圆的乘积，运动只剩两个相位随时间变化。频率比为有理数时轨道闭合，频率比为无理数时轨道在环面上稠密。这个例子提出两个不同问题：环面是否来自振子方程的特殊形式？即使每条纤维都是环面，能否在所有纤维上同时选择一致的相位和作用？

\ProseParagraphBreak

第九章已经说明，Poisson 对合的函数产生交换的 Hamilton 场。若在 \(2n\) 维相空间中有 \(n\) 个独立的这类积分，它们的共同水平集恰有 \(n\) 维，而这些交换场张成全部切空间。紧性进一步使流完备，返回同一点的时间组成满秩格。这就从动力学本身构造出了环面。

\ProseParagraphBreak

邻域中的辛坐标还需要更多工作：相邻环面的周期必须能光滑选择，周期所确定的闭一形式必须能积分，角坐标中残留的底空间二形式必须能消去。本章逐步完成这些构造，再检查它们沿大闭路延拓时会发生什么。矩映射的凸性和环面作用的分类将在下一章展开；这里先把可积系统的局部坐标及其整体障碍说清楚。

\ProseParagraphBreak

本章沿用 \(\iota_{X_H}\omega=\mathrm{d}H\)、\(\omega=-\mathrm{d}\lambda\) 的符号，角变量统一取周期 \(2\pi\)。主要先修为第八章的 Lagrange 邻域定理、第九章的 Poisson 对合，以及第一卷的流与离散格。可积系统的基本定义、环面结构及球摆模型参考 Cannas da Silva 第5.5节；其中邻域坐标的证明由本章完整展开。\cite[第150--153页]{Cannas2006Symplectic}

% !TeX root = ../../Book4.tex
% 纲要标题：### 11.1 可积性与交换流
% 纲要位置：几何学计划纲要.md，第 3885 行。
% 正文状态：初稿；尚未正式审核。

% 纲要内容开始
% 
% * 独立第一积分及 Poisson 交换
% * 共同水平集的 Lagrange 性
% * 完备交换流与轨道
% * 紧连通正则纤维的环面识别
% 
% 纲要内容结束

\section{可积性与交换流}
\label{g4:sec:1101}

一个第一积分把运动限制在一个水平集上。若在 \(2n\) 维相空间中找到 \(n\) 个独立且彼此 Poisson 对合的第一积分，共同水平集还会带有 \(n\) 个相互交换的运动方向。这些方向恰好铺满正则纤维的切空间。紧性使运动可以持续任意时间，而重复回到原点的时间向量组成一个格；环面结构便从这个格产生。本节逐步建立这条联系，下一节再把一张纤维上的结构延伸到它的邻域。

\subsection{独立第一积分及 Poisson 交换}
\label{g4:subsec:1101:01}

以下 \((M,\omega)\) 为无边界辛流形，实维数为 \(2n\)，其中 \(n\geq1\)。沿用第九章的约定
\[
	\iota_{X_f}\omega=\mathrm{d}f,\qquad
	\{f,g\}=\omega(X_f,X_g)=X_gf.
\]
因此 \(f\) 是自主 Hamilton 系统 \(H\) 的第一积分，当且仅当 \(\{f,H\}=0\)。这使 \(f\) 沿 \(H\) 的轨道保持常值；若要把几个第一积分产生的运动同时组织起来，还要要求它们彼此对合。

\begin{definition}\label{g4:def:1101:integrable-system}
	设 \(f_1,\ldots,f_n\in C^\infty(M,\mathbb R)\)，记
	\[
		F=(f_1,\ldots,f_n):M\longrightarrow\mathbb R^n.
	\]
	若 \(\mathrm{d}f_1(x),\ldots,\mathrm{d}f_n(x)\) 线性无关，就说这些函数在 \(x\) 处\term[hanshuduli]{函数独立}[functionally independent]。若
	\begin{equation}\label{g4:eq:1101:integrability}
		\{f_i,f_j\}=0\quad(1\leq i,j\leq n),\qquad
		M_{\mathrm{reg}}\coloneqq\{x\in M:\operatorname{rank}\mathrm{d}F_x=n\}
		\text{ 在 }M\text{ 中稠密},
	\end{equation}
	则称 \((M,\omega,F)\) 为\term[wanquankejixitong]{完全可积系统}[completely integrable system]，也简称可积系统。对一个指定的 Hamilton 函数 \(H\)，若可取这样的积分组并使 \(f_1=H\)，则称 \(H\) 完全可积。
\end{definition}

函数独立是一阶条件。例如 \(f\) 与 \(f^2\) 即使作为函数不成常数倍，它们的微分仍满足 \(\mathrm{d}(f^2)=2f\,\mathrm{d}f\)，不能提供两个独立约束。满秩条件在小扰动下保持，所以 \(M_{\mathrm{reg}}\) 是开集；稠密性允许系统在较小的集合上出现奇异纤维。定义中的独立性和对合条件与通常的完全可积系统定义一致。\cite[第5.5节，定义5.13，第151页]{Cannas2006Symplectic}

\ProseParagraphBreak

这里为什么恰好取 \(n\) 个函数？在一点 \(x\)，非退化性使
\[
	\mathrm{T}_xM\longrightarrow\mathrm{T}_x^*M,
	\qquad v\longmapsto\iota_v\omega_x
\]
为线性同构。因此 \(k\) 个函数的微分独立，当且仅当它们的 Hamilton 向量在该点独立；若函数又两两对合，这些向量的张成空间满足
\[
	\omega_x(X_{f_i}(x),X_{f_j}(x))=\{f_i,f_j\}(x)=0.
\]
这是一个各向同性子空间。由\cref{g4:prop:0701:lagrangian-characterizations}及其前面的维数结论，它的维数至多为 \(n\)，故 \(k\leq n\)。完全可积性正好达到独立对合函数数目的这一上限。

\ProseParagraphBreak

若 \(h\) 是 \(F(M)\) 的一个开邻域上的光滑函数，复合 Hamilton 函数 \(H=h\circ F\) 也保持全部 \(f_i\)。更具体地，链式法则和辛形式的非退化性给出
\begin{equation}\label{g4:eq:1101:composite-hamiltonian}
	\mathrm{d}H=\sum_{i=1}^n
	\frac{\partial h}{\partial c_i}(F)\,\mathrm{d}f_i,
	\qquad
	X_H=\sum_{i=1}^n
	\frac{\partial h}{\partial c_i}(F)X_{f_i}.
\end{equation}
这些系数在每个共同水平集上都是常数。于是同一组积分还可以描述许多不同的 Hamilton 运动；即使某处 \(\mathrm{d}H=0\)，积分映射 \(F\) 仍可能正则。

\begin{example}\label{g4:exm:1101:momentum-fibres}
	在 \(\mathbb R^{2n}_{q,p}\) 上取 \(\omega=\sum_i\mathrm{d}q_i\wedge\mathrm{d}p_i\) 和 \(f_i=p_i\)。微分 \(\mathrm{d}p_i\) 处处独立，且 \(\{p_i,p_j\}=0\)。相应向量场及流为
	\[
		X_{f_i}=\partial_{q_i},\qquad
		\varphi_i^t(q,p)=(q+t e_i,p).
	\]
	每张纤维 \(F^{-1}(c)=\mathbb R^n_q\times\{c\}\) 都由这些平移方向铺满。对 \(H=\frac12\sum_i p_i^2\)，式\eqref{g4:eq:1101:composite-hamiltonian}给出 \(\dot q=c\)、\(\dot p=0\)。当 \(c=0\) 时，\(H\) 的运动静止，但 \(F\) 仍处处满秩；单个 Hamilton 函数的临界性与积分映射的奇异性是两回事。
\end{example}

\subsection{共同水平集的 Lagrange 性}
\label{g4:subsec:1101:02}

给定 \(c\in\mathbb R^n\)，我们将考察 \(F^{-1}(c)\) 的一个非空连通分支 \(L\)。若 \(\operatorname{rank}\mathrm{d}F_x=n\) 对每个 \(x\in L\) 成立，就称 \(L\) 为一张\term[zhengzeliantongxianwei]{正则连通纤维}[regular connected fibre]。当 \(c\) 是正则值时，它的每个非空连通纤维都满足这一条件；这个定义也允许同一水平集的其他分支含有临界点。后面所说的紧纤维，都是指选定的 \(L\) 紧。

\begin{proposition}\label{g4:prop:1101:regular-lagrangian-level}
	设 \((M,\omega)\) 为实维数 \(2n\) 的无边界辛流形，\(F=(f_1,\ldots,f_n):M\to\mathbb R^n\) 的各分量两两 Poisson 对合。设 \(L\) 是 \(F^{-1}(c)\) 的一个非空连通分支，并且 \(\operatorname{rank}\mathrm{d}F_x=n\) 对所有 \(x\in L\) 成立。则 \(L\) 是无边界的嵌入 Lagrange 子流形，且
	\begin{equation}\label{g4:eq:1101:tangent-frame}
		\mathrm{T}_xL=\ker\mathrm{d}F_x
		=\operatorname{span}_{\mathbb R}
		\{X_{f_1}(x),\ldots,X_{f_n}(x)\}.
	\end{equation}
	这些 Hamilton 向量场限制到 \(L\) 后构成处处线性无关、两两交换的光滑向量场。
\end{proposition}
\begin{proof}
	先确认 \(L\) 的流形结构。对 \(x\in L\)，满秩条件在某个开邻域内保持。第一卷定理11.4.4的正则水平集定理应用于该邻域，给出局部坐标，其中 \(F=c\) 固定前 \(n\) 个坐标，其余 \(n\) 个坐标自由变化。把坐标域缩为小盒，可使这个局部水平集连通；它含有 \(x\)，所以全在 \(x\) 所属的分支 \(L\) 内。这些坐标因此给出 \(L\) 的嵌入子流形图，并且
	\[
		\dim L=n,\qquad
		\mathrm{T}_xL=\bigcap_{j=1}^n\ker\mathrm{d}f_j(x).
	\]
	这个论证只在 \(L\) 附近使用满秩性。

	对任意 \(i,j\)，有
	\[
		\mathrm{d}f_j(X_{f_i})=\{f_j,f_i\}=0,
	\]
	所以 \(X_{f_i}\) 切于 \(L\)。若 \(\sum_i a_iX_{f_i}(x)=0\)，与 \(\omega_x\) 收缩后得到 \(\sum_i a_i\mathrm{d}f_i(x)=0\)，故全部 \(a_i=0\)。这 \(n\) 个独立切向量构成 \(\mathrm{T}_xL\) 的基，证明式\eqref{g4:eq:1101:tangent-frame}。它们两两的辛配对都为零，因而 \(\omega|_{\mathrm{T}_xL}=0\)；半维各向同性判别说明 \(L\) 是 Lagrange 子流形。

	最后，\cref{g4:prop:0902:hamiltonian-bracket}给出
	\[
		[X_{f_i},X_{f_j}]=-X_{\{f_i,f_j\}}=0.
	\]
	切于子流形的光滑向量场可限制成该子流形上的光滑向量场，其括号也限制为原括号。例如在适配坐标中，法向分量在子流形上为零，切向导数仍为零；切向分量的括号公式正是受限向量场的括号。因此受限向量场同样两两交换。
\end{proof}

共同水平集的维数与其上的运动方向数在这里相等。只有“沿原运动守恒”还不能保证这一点：若第一积分之间的括号不为零，它们产生的运动可能离开彼此的水平集；对合条件正是使全部运动同时切于 \(L\) 的条件。

\begin{example}\label{g4:exm:1101:noncommuting-level}
	在 \(\mathbb R^4_{q_1,q_2,p_1,p_2}\) 上，函数 \(q_1,p_1\) 处处独立，故共同水平集 \(q_1=a,p_1=b\) 是二维平面。但 \(\{q_1,p_1\}=1\)，而该平面上的限制辛形式为 \(\mathrm{d}q_2\wedge\mathrm{d}p_2\)，所以它是辛平面。相应向量场 \(X_{q_1}=-\partial_{p_1}\)、\(X_{p_1}=\partial_{q_1}\) 虽然交换，却都不切于这个共同水平集。这个例子同时说明，对合条件要求函数括号为零，单独要求 Hamilton 向量场交换还不够。
\end{example}

\subsection{完备交换流与轨道}
\label{g4:subsec:1101:03}

现在在正则连通纤维 \(L\) 上记 \(Y_i=X_{f_i}|_L\)。上一命题给出的是一组交换标架。要把任意时间向量 \(t=(t_1,\ldots,t_n)\) 解释为运动，就要使每个 \(Y_i\) 的积分曲线对全部实数时间存在；这正是完备性。这里讨论的向量场定义在 \(L\) 上。

\begin{proposition}\label{g4:prop:1101:complete-action}
	设 \(L\) 为非空连通的 \(n\) 维无边界光滑流形，\(Y_1,\ldots,Y_n\) 为处处线性无关、两两交换的完备光滑向量场。记它们的全局流为 \(\varphi_i^t\)。则
	\begin{equation}\label{g4:eq:1101:joint-flow}
		\Phi:\mathbb R^n\times L\longrightarrow L,
		\qquad
		\Phi(t,x)=\varphi_1^{t_1}\circ\cdots\circ\varphi_n^{t_n}(x)
	\end{equation}
	是光滑传递的加法群作用。对任意 \(x_0\in L\)，子群
	\begin{equation}\label{g4:eq:1101:period-subgroup}
		\Gamma=\{t\in\mathbb R^n:\Phi(t,x_0)=x_0\}
	\end{equation}
	闭且离散，与所选 \(x_0\) 无关。轨道映射诱导微分同胚
	\begin{equation}\label{g4:eq:1101:homogeneous-fibre}
		\overline P_{x_0}:\mathbb R^n/\Gamma\longrightarrow L,
		\qquad [t]\longmapsto\Phi(t,x_0).
	\end{equation}
\end{proposition}
\begin{proof}
	第一卷定理13.1.7说明，每个完备场的流是光滑映射 \(\mathbb R\times L\to L\)，且 \(\varphi_i^{-t}\) 为 \(\varphi_i^t\) 的逆。先把两两流的交换性落实到任意实数时间。第一卷命题13.3.3的拉回求导公式与 \([Y_i,Y_j]=0\) 给出
	\[
		\frac{\mathrm{d}}{\mathrm{d}s}(\varphi_i^s)^*Y_j=0,
		\qquad
		\mathrm{d}\varphi_i^s(Y_j)=Y_j\circ\varphi_i^s.
	\]
	固定任意 \(s\in\mathbb R\) 后，曲线 \(t\mapsto\varphi_i^s(\varphi_j^t(x))\) 是 \(Y_j\) 的积分曲线，初值为 \(\varphi_i^s(x)\)。它对全部 \(t\) 有定义，解的唯一性于是给出
	\[
		\varphi_i^s\circ\varphi_j^t
		=\varphi_j^t\circ\varphi_i^s
		\qquad(s,t\in\mathbb R).
	\]
	所以式\eqref{g4:eq:1101:joint-flow}的复合次序无关，并满足 \(\Phi(0,x)=x\) 及 \(\Phi(s,\Phi(t,x))=\Phi(s+t,x)\)，这就是所需群作用。

	对每个 \(x\)，令 \(P_x(t)=\Phi(t,x)\)。在零点求微分，得到
	\begin{equation}\label{g4:eq:1101:orbit-differential}
		\mathrm{d}(P_x)_0(a_1,\ldots,a_n)
		=\sum_{i=1}^n a_iY_i(x).
	\end{equation}
	这由标架假设是线性同构。由第一卷定理11.3.3，\(P_x\) 把零点的一个邻域微分同胚地映到 \(x\) 的一个开邻域，所以每条轨道都在其每一点附近包含开邻域，即每条轨道开。轨道彼此相同或不相交；若有两条以上轨道，一条轨道与其余轨道的并就把 \(L\) 分成两个非空开集，违背连通性。因此作用传递。群作用公式还给出
	\[
		P_{x_0}(t+a)=P_{P_{x_0}(t)}(a),
	\]
	所以 \(P_{x_0}\) 在每个 \(t\) 处都是局部微分同胚。

	子群 \(\Gamma=P_{x_0}^{-1}(\{x_0\})\) 闭；\(P_{x_0}\) 在零点附近单射，所以零点是 \(\Gamma\) 的孤立点，再由平移可知每个点都孤立。若 \(y=\Phi(s,x_0)\)，加法的交换性给出
	\[
		\Phi(t,y)=y
		\quad\Longleftrightarrow\quad
		\Phi(-s,\Phi(t,\Phi(s,x_0)))=x_0
		\quad\Longleftrightarrow\quad
		\Phi(t,x_0)=x_0.
	\]
	传递性说明全部点都有同一稳定子。

	第一卷定理17.1.5应用于加法 Lie 群 \(\mathbb R^n\) 的闭子群 \(\Gamma\)，赋予商 \(\mathbb R^n/\Gamma\) 光滑齐性空间结构。这里还可直接看见局部商图：选一个零点邻域 \(V\)，使 \((V-V)\cap\Gamma=\{0\}\)，则商映射在每个 \(t+V\) 上单射，给出局部微分同胚。两点 \(s,t\) 有相同轨道像，当且仅当 \(s-t\in\Gamma\)，故式\eqref{g4:eq:1101:homogeneous-fibre}良定义且为双射。其局部表达由 \(P_{x_0}\) 与上述商图给出，是局部微分同胚。双射的局部光滑逆在交集上一致，便拼成全局光滑逆。
\end{proof}

选择 \(x_0\) 将 \(L\) 与群 \(\mathbb R^n/\Gamma\) 识别，并把 \(x_0\) 选为单位元；换一个基点会改变这一识别的平移原点。子群 \(\Gamma\) 则由给定的交换标架决定。若 \(\Gamma\neq\{0\}\)，\(\mathbb R^n\) 的作用不是自由作用；除去这些共同周期后，\(\mathbb R^n/\Gamma\) 在 \(L\) 上自由且传递。

\begin{example}\label{g4:exm:1101:incomplete-fibre}
	取 \(M=(0,1)_q\times\mathbb R_p\)、\(\omega=\mathrm{d}q\wedge\mathrm{d}p\) 和 \(f=p\)。每张纤维 \(L_c=(0,1)\times\{c\}\) 都正则，且 \(X_f=\partial_q\) 在其切空间中处处非零。过 \((q_0,c)\) 的积分曲线为 \((q_0+t,c)\)，最大时间区间是 \((-q_0,1-q_0)\)。局部流因此存在且保持纤维，但不能用全部实数时间定义式\eqref{g4:eq:1101:joint-flow}。完备性所补足的正是从局部运动到群作用的这一步。
\end{example}

\subsection{紧连通正则纤维的环面识别}
\label{g4:subsec:1101:04}

离散子群 \(\Gamma\) 决定哪些时间方向会闭合。第一卷引理19.2.2证明：任意 \(\mathbb R^n\) 的离散加法子群都有一组实线性无关的整数基，即存在 \(0\leq r\leq n\) 和实线性无关的 \(\gamma_1,\ldots,\gamma_r\)，使
\begin{equation}\label{g4:eq:1101:lattice-basis}
	\Gamma=\mathbb Z\gamma_1\oplus\cdots\oplus\mathbb Z\gamma_r.
\end{equation}
其证明先取最短非零群元素，再投影到其正交补并归纳。这里投影仍离散的原因是：沿最短元素减去整数倍后，纵向分量落在一个有界线段内；若非零投影趋于零，就会得到有界区域内无限多个群元素，与离散子群的点间距有正下界矛盾。这一论证保证了整数基的存在，避免把实线性基误当成整数生成元。

\begin{proposition}\label{g4:prop:1101:complete-fibre-cylinder}
	在\cref{g4:prop:1101:regular-lagrangian-level}的假设下，若 \(X_{f_1}|_L,\ldots,X_{f_n}|_L\) 全部完备，则存在唯一整数 \(r\in\{0,\ldots,n\}\)，使其周期子群的秩为 \(r\)，且有微分同胚
	\[
		L\cong\mathbb R^{n-r}\times\mathbb T^r,
		\qquad \mathbb T^r\coloneqq
		(\mathbb R/(2\pi\mathbb Z))^r.
	\]
	特别地，\(L\) 紧当且仅当 \(r=n\)。
\end{proposition}
\begin{proof}
	由\cref{g4:prop:1101:complete-action}，\(L\cong\mathbb R^n/\Gamma\)。取式\eqref{g4:eq:1101:lattice-basis}的整数基，并补成 \(\mathbb R^n\) 的实线性基 \(\gamma_1,\ldots,\gamma_n\)。若 \(t=\sum_{j=1}^n a_j\gamma_j\)，则
	\[
		[t]\longmapsto
		\bigl(a_{r+1},\ldots,a_n,
		2\pi a_1\bmod2\pi,\ldots,2\pi a_r\bmod2\pi\bigr)
	\]
	是所需微分同胚：加上 \(\Gamma\) 的元素恰好将前 \(r\) 个系数加上整数，后 \(n-r\) 个系数不变，故映射良定义且双射；在线性坐标与商的局部图中，它及其逆都光滑。整数 \(r\) 等于 \(\dim_{\mathbb R}\operatorname{span}_{\mathbb R}\Gamma\)，因而由 \(\Gamma\) 唯一决定。

	当 \(r=n\) 时，紧立方体 \([0,2\pi]^n\) 满射到 \(\mathbb T^n\)，所以 \(L\) 紧。当 \(r<n\) 时，向 \(\mathbb R^{n-r}\) 的投影是连续满射；若 \(L\) 紧，像也应紧，这与正维 Euclid 空间不紧矛盾。因此紧性恰好迫使全部方向都成为周期方向。
\end{proof}

下面的结论称为\term[liouvillehuanmiandingli]{Liouville 环面定理}[Liouville torus theorem]。它把纤维的紧性转化为完整周期格的满秩性。

\begin{theorem}[Liouville 环面定理]\label{g4:thm:1101:liouville-torus}
	设 \((M,\omega)\) 为实维数 \(2n\) 的无边界辛流形，\(F=(f_1,\ldots,f_n):M\to\mathbb R^n\) 的各分量两两 Poisson 对合。若 \(L\) 是 \(F^{-1}(c)\) 的一个非空紧连通分支，且 \(\operatorname{rank}\mathrm{d}F_x=n\) 在 \(L\) 上处处成立，则 \(L\) 是 Lagrange 子流形，并微分同胚于 \(\mathbb T^n\)。更具体地，受限 Hamilton 向量场产生传递的 \(\mathbb R^n\) 作用，其共同周期子群是满秩格
	\[
		\Gamma=\mathbb Z\gamma_1\oplus\cdots\oplus\mathbb Z\gamma_n.
	\]
	选定 \(x_0\in L\) 后，下式给出一个微分同胚：
	\begin{equation}\label{g4:eq:1101:torus-parametrization}
		\chi:\mathbb T^n\longrightarrow L,
		\qquad
		\chi(\theta)=\Phi\left(\frac1{2\pi}
		\sum_{j=1}^n\theta_j\gamma_j,x_0\right).
	\end{equation}
\end{theorem}
\begin{proof}
	\cref{g4:prop:1101:regular-lagrangian-level}说明 \(L\) 是无边界的 \(n\) 维 Lagrange 子流形，并给出交换标架 \(Y_i=X_{f_i}|_L\)。第一卷定理13.2.5说明，紧无边界流形上的每个光滑向量场都完备。该结果来自有限时间逃逸判据：若最大积分曲线在有限时间终止，就必须离开每个紧集，而这里整个 \(L\) 已经紧。因此每个 \(Y_i\) 都完备。这些在 \(L\) 上的全局曲线作为 \(M\) 中的曲线也解原方程，唯一性使它们与从 \(L\) 出发的环境积分曲线一致。

	现在可用\cref{g4:prop:1101:complete-action}得到 \(L\cong\mathbb R^n/\Gamma\)，再由\cref{g4:prop:1101:complete-fibre-cylinder}和紧性得 \(\operatorname{rank}\Gamma=n\)。式\eqref{g4:eq:1101:torus-parametrization}是该商同构在格基下的坐标表达。把某个 \(\theta_j\) 加上 \(2\pi\) 恰好使时间向量加上 \(\gamma_j\)，故表达式良定义；格基的整数独立性给出单射性，实线性基性质与作用的传递性给出满射性，局部商图给出双向光滑性。\cite[第5.5节，定理5.14(a)，第152页]{Cannas2006Symplectic}
\end{proof}

称这样的 \(L\) 为\term[liouvillehuanmian]{Liouville 环面}[Liouville torus]，其满秩格 \(\Gamma\) 称为该积分组在 \(L\) 上的\term[zhouqige]{周期格}[period lattice]。紧性只用于所选分支；证明中没有要求整个水平集紧，也没有要求环境中其他初值的 Hamilton 流完备。

\ProseParagraphBreak

周期格一般不由各个 \(X_{f_i}\) 单独的周期产生：某个 \(X_{f_i}\) 可能没有非零周期，若干流的适当组合却会闭合。写
\[
	B=(\gamma_1\ \cdots\ \gamma_n),
	\qquad
	\gamma_j=(\gamma_j^1,\ldots,\gamma_j^n)^{\mathsf T},
\]
则式\eqref{g4:eq:1101:torus-parametrization}给出
\begin{equation}\label{g4:eq:1101:normalized-fields}
	\chi_*\partial_{\theta_j}
	=\frac1{2\pi}\sum_{i=1}^n\gamma_j^iX_{f_i}|_L,
	\qquad
	X_{f_i}|_L
	=\chi_*\left(2\pi\sum_{j=1}^n
	(B^{-1})_{ji}\partial_{\theta_j}\right).
\end{equation}
第一式中的场以 \(2\pi\) 为最小正周期，因为 \(\gamma_j\) 是整数基中的一个元素。对 \(H=h\circ F\)，结合式\eqref{g4:eq:1101:composite-hamiltonian}便得到在这张环面上的运动
\begin{equation}\label{g4:eq:1101:torus-frequency}
	\theta(t)=\theta(0)+t\nu\pmod{2\pi\mathbb Z^n},
	\qquad \nu=2\pi B^{-1}\nabla h(c).
\end{equation}
速度是常向量，所以原来的 Hamilton 运动在这些角参数中成为匀速平移。

\begin{example}\label{g4:exm:1101:cylinder-periods}
	设 \(Q=\mathbb T^r\times\mathbb R^{n-r}\)，在 \(\mathrm{T}^*Q\) 上用 \((q_1,\ldots,q_n,p_1,\ldots,p_n)\) 表示坐标，前 \(r\) 个 \(q_j\) 按模 \(2\pi\) 理解。形式 \(\omega=\sum_i\mathrm{d}q_i\wedge\mathrm{d}p_i\) 全局有定义，取 \(F=(p_1,\ldots,p_n)\)。在每张纤维上，联合流为 \(q\mapsto q+t\)，周期子群恰为
	\[
		\Gamma=(2\pi\mathbb Z)^r\times\{0\}^{n-r}.
	\]
	因此 \(r\) 个闭合方向产生圆周，其余方向保持为实直线；所有向量场都完备，而纤维只有在 \(r=n\) 时紧。这给出了上一分类中每一种可能的秩。
\end{example}

本节构造的角参数定义在一张固定的纤维上，周期格也暂时固定在这张纤维。为了比较相邻环面上的运动，需要让周期基随纤维光滑变化，并找到与这些角参数配对的横向坐标。下一节用作用积分完成这一构造，得到邻域上的辛形式 \(\sum_i\mathrm{d}\theta_i\wedge\mathrm{d}I_i\)。
% 纲要条目排为编号小节标题，并进入目录与书签；说明段仍保留于纲要注释中。

% !TeX root = ../../Book4.tex
% 纲要标题：### 11.2 Liouville–Arnold 定理
% 纲要位置：几何学计划纲要.md，第 3892 行。
% 正文状态：初稿；尚未正式审核。

% 纲要内容开始
% 
% * 周期格和局部基底
% * 作用积分的定义
% * 角坐标的构造
% * 辛形式的标准表达
% * 参数及格基变化
% 
% 纲要内容结束

\section{Liouville–Arnold 定理}
\label{g4:sec:1102}

本节证明\term[liouvillearnolddingli]{Liouville--Arnold 定理}[Liouville--Arnold theorem]。\footnote{\OriginalName{Владимир Игоревич Арнольд}（Vladimir Igorevich Arnold，1937--2010），苏联、俄罗斯数学家，参与建立 Hamilton 动力学的 KAM 理论，并研究奇点理论。}一个紧正则共同水平集给出不变环面，但邻近轨道怎样分布、每个环面的周期怎样变化，还需要在环面以外选择坐标。作用角定理将这些问题合在一起：先把紧环面扩展为一个光滑环面族，再把周期归一化为 \(2\pi\)，最后选择与这些角变量配对的作用变量。作用变量由沿闭路的一形式积分得到，通常不是最初给定的第一积分本身。

\subsection{周期格和局部基底}
\label{g4:subsec:1102:01}

本节只研究一个指定紧组件附近的几何。它允许同一共同水平集在远处还有其他组件，也不要求第一积分组成的映射在整个流形上固有。首先构造一个足够小的开邻域，使所选组件及其邻近纤维都完整地包含在其中。

\begin{lemma}\label{g4:lem:1102:compact-fibration}
	设 \((M^{2n},\omega)\)、\(n\ge1\)，为无边界、Hausdorff 且第二可数的辛流形，\(F=(f_1,\ldots,f_n):M\to\mathbb R^n\) 光滑，且 \(\{f_i,f_j\}=0\)。设 \(L\) 为 \(F^{-1}(b_0)\) 的非空紧连通分支，\(\operatorname{rank}\mathrm{d}F=n\) 在 \(L\) 上处处成立。则存在 \(b_0\) 的开球邻域 \(B\)、\(L\) 的开邻域 \(U\subset M\)，以及微分同胚
	\[
		\chi:L\times B\longrightarrow U,\qquad
		F(\chi(q,b))=b,\qquad \chi(q,b_0)=q.
	\]
	每个 \(L_b=\chi(L\times\{b\})\) 都是 \(F^{-1}(b)\) 的紧连通分支，并为 Lagrange 环面；\(F|_U:U\to B\) 是固有淹没。此外，存在 \(\lambda\in\Omega^1(U)\)，使 \(\omega|_U=-\mathrm{d}\lambda\) 且 \(\lambda|_{\mathrm{T}L}=0\)。
\end{lemma}
\begin{proof}
	由\cref{g4:thm:1101:liouville-torus}，\(L\) 是嵌入的 Lagrange 环面。它紧，因而在 Hausdorff 流形 \(M\) 中闭。\cref{g4:thm:0803:lagrangian-neighborhood}给出零截面附近的辛同胚
	\[
		\Psi:V\subset\mathrm{T}^*L\longrightarrow W\subset M,\qquad
		\Psi(q,0)=q,\quad \Psi^*\omega=-\mathrm{d}\theta_L.
	\]
	记 \(\widetilde F=F\Psi\)。它沿零截面恒为 \(b_0\)，所以其微分在零截面的切向上为零。总秩为 \(n\)，故竖直微分
	\[
		A_q=\mathrm{d}_p\widetilde F(q,0):
		\mathrm{T}_q^*L\longrightarrow\mathbb R^n
	\]
	是线性同构，且依赖 \(q\) 光滑。

	用 \(v=A_qp\) 作纤维变量，定义
	\(G(q,v)=\widetilde F(q,A_q^{-1}v)-b_0\)。有
	\(G(q,0)=0\)、\(\mathrm{D}_vG(q,0)=I\)。紧性使我们能选同一个 \(r>0\)，使所有 \((q,A_q^{-1}v)\)、\(|v|\le r\)，均属于 \(V\)，并且
	\[
		\sup_{q\in L,\ |v|\le r}\|\mathrm{D}_vG(q,v)-I\|\le\frac12.
	\]
	第一卷定理11.3.2的定量局部逆说明，对每个 \(q\) 和 \(|b-b_0|<r/2\)，方程 \(G(q,v)=b-b_0\) 在 \(|v|<r\) 内有唯一解。第一卷定理11.3.5的参数隐函数定理说明该解联合依赖 \((q,b)\) 光滑：在任一 \(q\) 坐标域中应用该定理，所得局部解由刚才的唯一性相容。因此可写成 \(p=P(q,b)\)。

	取 \(B=B_{r/2}(b_0)\)，令 \(\chi(q,b)=\Psi(q,P(q,b))\)。映射
	\((q,p)\mapsto(q,\widetilde F(q,p))\) 的竖直导数可逆，所以 \(\chi\) 是局部微分同胚；它的全局单射性来自保留 \(q\) 及上述纤维解的唯一性。于是其像 \(U\) 开，\(\chi:L\times B\to U\) 是微分同胚。

	在 \(U\) 中，\(F\) 是投影，故满秩，且 \(L_b=F^{-1}(b)\cap U\) 紧连通。这个集合在整个 \(F^{-1}(b)\) 中开；它在 \(M\) 中紧且闭，故在该水平集中也闭。因此它确为一个完整连通分支。每个 \(L_b\) 再由\cref{g4:thm:1101:liouville-torus}成为 Lagrange 环面。对紧集 \(K\subset B\)，\((F|_U)^{-1}(K)=\chi(L\times K)\) 紧，所以这里的限制映射固有。最后在 \(U\) 上取
	\(\lambda=(\Psi^{-1})^*\theta_L\)，便有所需的原始形式，并且它在原环面的切向上为零。
\end{proof}

这个引理中的 \(U\) 是由完整纤维组件组成的邻域，不必等于整个 \(F^{-1}(B)\)。它已经足以使各个 \(X_{f_j}\) 在 \(U\) 上完备：从任一点出发的积分曲线始终留在一个紧 \(L_b\) 中，有限时间不能逃出。这些完备场的光滑流依次复合，利用流的交换性给出光滑作用
\[
	\Phi:\mathbb R^n\times U\longrightarrow U,\qquad
	\Phi(t,x)=\Phi^{X_{f_1}}_{t_1}\cdots\Phi^{X_{f_n}}_{t_n}(x).
\]
由\cref{g4:prop:1101:complete-action}在各纤维上的结论，每个 \(L_b\) 都是一条轨道。取光滑截面 \(s(b)=\chi(q_0,b)\)，并令
\begin{equation}\label{g4:eq:1102:period-lattice}
	\Gamma_b=\{t\in\mathbb R^n:\Phi(t,s(b))=s(b)\}.
\end{equation}
因为作用交换且传递，这个稳定子不依赖纤维中的起点。上一节说明 \(\Gamma_b\) 为满格；现在还要控制它随 \(b\) 的变化。

\begin{lemma}\label{g4:lem:1102:smooth-period-basis}
	在\cref{g4:lem:1102:compact-fibration}构造的环面族中，任取 \(\Gamma_{b_0}\) 的一个整数基底 \(\ell_1,\ldots,\ell_n\)。缩小 \(B\) 后，存在光滑函数
	\(\lambda_i:B\to\mathbb R^n\)，使
	\[
		\lambda_i(b_0)=\ell_i,\qquad
		\Gamma_b=\Lambda(b)\mathbb Z^n,\qquad
		\Lambda(b)=\bigl(\lambda_1(b)\ \cdots\ \lambda_n(b)\bigr).
	\]
	特别地，延拓出的 \(n\) 个周期构成完整周期格的基底。
\end{lemma}
\begin{proof}
	令 \(E(t,b)=\Phi(t,s(b))\)。在任意 \((t,b)\) 处，\(t\) 方向的微分由 \(n\) 个独立场 \(X_{f_j}\) 给出，张成纤维切空间；而 \(F(E(t,b))=b\)，故 \(b\) 方向的微分投影为恒等映射。因此 \(E\) 处处为局部微分同胚。

	在 \((\ell_i,b_0)\) 附近使用其局部逆，并把目标取成 \(s(b)\)，就得到唯一光滑的 \(\lambda_i(b)\)，满足 \(E(\lambda_i(b),b)=s(b)\)。由于 \(\Lambda(b_0)\) 可逆，缩小 \(B\) 后 \(\Lambda(b)\) 仍可逆。于是
	\(\Gamma'_b=\Lambda(b)\mathbb Z^n\) 是 \(\Gamma_b\) 的子格。还需证明没有遗漏其他周期。

	先在 \((0,b_0)\) 处应用同一个局部逆。存在 \(\varepsilon>0\) 和共同的参数邻域，使
	\[
		|t|<\varepsilon,\quad E(t,b)=s(b)
		\quad\Longrightarrow\quad t=0.
	\]
	这是对邻近所有纤维都成立的无短周期估计。若任意小参数邻域内都有 \(\Gamma'_b\ne\Gamma_b\)，可取 \(b_\nu\to b_0\) 及
	\(t_\nu\in\Gamma_{b_\nu}\setminus\Gamma'_{b_\nu}\)。从 \(t_\nu\) 减去适当的 \(\Lambda(b_\nu)\mathbb Z^n\) 元素，得到
	\[
		r_\nu=\Lambda(b_\nu)a_\nu\in\Gamma_{b_\nu},
		\qquad a_\nu\in[-1/2,1/2]^n,\qquad r_\nu\ne0.
	\]
	取子列使 \(a_\nu\to a\)。由 \(E\) 连续，
	\(\Lambda(b_0)a\in\Gamma_{b_0}=\Lambda(b_0)\mathbb Z^n\)。
	故 \(a\in\mathbb Z^n\cap[-1/2,1/2]^n=\{0\}\)，从而 \(r_\nu\to0\)。这与统一无短周期估计矛盾。因此缩小参数邻域后 \(\Gamma'_b=\Gamma_b\)。
\end{proof}

由一个周期向量 \(\lambda_i(b)\) 得到的闭路写成
\begin{equation}\label{g4:eq:1102:period-cycle}
	\gamma_i(t,b)=\Phi(t\lambda_i(b),s(b)),\qquad t\in\mathbb R/\mathbb Z.
\end{equation}
这里参数走一周的时间是 \(1\)，真正的原 Hamilton 多时间增量是列向量 \(\lambda_i(b)\)。这一区分将决定作用积分中的 \(2\pi\) 因子。

\subsection{作用积分的定义}
\label{g4:subsec:1102:02}

前面的邻域已经有 \(\omega=-\mathrm{d}\lambda\)。在每个纤维上，\(\lambda\) 的限制闭，因为纤维是 Lagrange 子流形。因此它沿闭路的积分只取决于闭路的整数绕行类：同伦移动闭路时，微分为零；闭路拼接使积分相加，反向使积分变号。

\begin{definition}\label{g4:def:1102:action-variables}
	在上述邻域中，固定光滑周期基底及相应闭路族 \(\gamma_i\)。定义\term[zuoyongjifen]{作用积分}[action integral]与相应的\term[zuoyongbianliang]{作用变量}[action variable]为
	\begin{equation}\label{g4:eq:1102:action-integral}
		I_i(b)=\frac1{2\pi}\int_{\gamma_i(\cdot,b)}\lambda,
		\qquad i=1,\ldots,n.
	\end{equation}
	以下也用 \(I_i\) 表示 \(I_i\circ F\) 在 \(U\) 上的拉回。
\end{definition}

这个定义把闭路扫过的有向辛面积转成参数函数。关键在于它们是否独立；答案由周期矩阵给出。

\begin{proposition}\label{g4:prop:1102:action-derivative}
	对\eqref{g4:eq:1102:action-integral}定义的函数，有
	\begin{equation}\label{g4:eq:1102:action-derivative}
		\mathrm{d}I_i=\frac1{2\pi}\sum_{j=1}^n\lambda_{ji}(b)\,\mathrm{d}b_j,
		\qquad
		\mathrm{D}I(b)=\frac1{2\pi}\Lambda(b)^{\mathsf T}.
	\end{equation}
	因此 \(I=(I_1,\ldots,I_n)\) 在 \(b_0\) 附近是一组光滑坐标。
\end{proposition}
\begin{proof}
	闭路族及 \(\lambda\) 光滑，紧参数区间上的积分说明 \(I_i\) 光滑。取参数曲线 \(b(u)\)，令
	\(\Gamma(u,t)=\gamma_i(t,b(u))\)，并记
	\(V=\partial_u\Gamma\)、\(T=\partial_t\Gamma\)。
	将拉回一形式写成
	\(\Gamma^*\lambda=a(u,t)\,\mathrm{d}t+c(u,t)\,\mathrm{d}u\)。
	外微分的坐标公式给出
	\[
		\partial_u a-\partial_t c
		=\mathrm{d}\lambda(V,T).
	\]
	因为 \(t\) 是圆周参数，\(c(u,1)=c(u,0)\)，积分时边界项消失。所以
	\[
		\frac{\mathrm{d}}{\mathrm{d}u}\int_{\gamma_i(\cdot,b(u))}\lambda
		=\int_0^1\mathrm{d}\lambda(V,T)\,\mathrm{d}t
		=\int_0^1\omega(T,V)\,\mathrm{d}t.
	\]
	这里最后的正号同时使用了 \(\mathrm{d}\lambda=-\omega\) 与二形式的反对称性。

	沿所取闭路，
	\[
		T=\sum_j\lambda_{ji}(b(u))X_{f_j},\qquad
		\mathrm{d}f_j(V)=b'_j(u).
	\]
	由 \(\iota_{X_{f_j}}\omega=\mathrm{d}f_j\)，被积式于是恒等于
	\(\sum_j\lambda_{ji}(b(u))b'_j(u)\)，与 \(t\) 无关。除以 \(2\pi\) 得到微分公式。\(\Lambda(b_0)\) 可逆，第一卷定理11.3.3的逆函数定理给出作用坐标。
\end{proof}

现在进一步缩小 \(B\)，使 \(I:B\to A\) 为微分同胚，其中 \(A\subset\mathbb R^n\) 是开球。仍用 \(U\) 表示相应较小邻域。以原参数 \(b\) 表示的周期一形式
\(\sum_j\lambda_{ji}\,\mathrm{d}b_j\) 因而都是闭的；这不是任意变化的格基自动满足的条件，而是辛形式闭性经作用积分体现出的约束。取前面由余切模型得到的 \(\lambda\) 时，\(I(b_0)=0\)；采用其他原始形式则可能得到不同的常数。

\subsection{角坐标的构造}
\label{g4:subsec:1102:03}

原场 \(X_{f_j}\) 的周期随 \(b\) 改变。作用积分的微分恰好给出把它们归一化为 \(2\pi\) 周期的组合。定义
\begin{equation}\label{g4:eq:1102:normalized-fields}
	Y_i=\frac1{2\pi}\sum_{j=1}^n\lambda_{ji}(F)\,X_{f_j}.
\end{equation}
由\eqref{g4:eq:1102:action-derivative}，
\[
	\iota_{Y_i}\omega
	=\frac1{2\pi}\sum_j\lambda_{ji}(F)\,\mathrm{d}f_j
	=\mathrm{d}I_i,
	\qquad Y_i=X_{I_i}.
\]
系数在每个共同水平集上恒定，故 \(Y_i\) 的流在该纤维上就是原作用的多时间平移 \(t\lambda_i(b)/(2\pi)\)。它的周期为 \(2\pi\)，而全部 \(Y_i\) 合起来给出自由、传递的
\(\mathbb T^n=(\mathbb R/2\pi\mathbb Z)^n\) 作用；自由性正是完整周期格等式 \(\Gamma_b=\Lambda(b)\mathbb Z^n\) 的用途。

\ProseParagraphBreak

保留初始截面 \(s\)，定义
\begin{equation}\label{g4:eq:1102:initial-angle-map}
	\Psi_0:\mathbb T^n\times A\longrightarrow U,\qquad
	\Psi_0(\vartheta,I)
	=\Phi\left(\frac1{2\pi}\Lambda(b(I))\vartheta,\ s(b(I))\right).
\end{equation}
改变 \(\vartheta\) 的代表只增加一个周期向量，所以定义良好。它在每个纤维上是自由传递作用的轨道参数化，故双射；其角方向微分为 \(Y_i\)，横向微分投影为 \(I\) 方向的恒等映射，所以是局部微分同胚。由此 \(\Psi_0\) 是整体微分同胚，给出了圆周值的\term[jiaobianliang]{角变量}[angle variable] \(\vartheta\)。

\begin{proposition}\label{g4:prop:1102:angle-form}
	在\eqref{g4:eq:1102:initial-angle-map}给出的坐标中，存在 \(A\) 上的光滑闭二形式 \(\beta\)，使
	\begin{equation}\label{g4:eq:1102:angle-base-form}
		\Psi_0^*\omega
		=\sum_{i=1}^n\mathrm{d}\vartheta_i\wedge\mathrm{d}I_i
		+\operatorname{pr}_A^*\beta.
	\end{equation}
	其中 \(\beta\) 正是初始截面在 \(I\) 参数下对 \(\omega\) 的拉回。
\end{proposition}
\begin{proof}
	记 \(\Omega=\Psi_0^*\omega\)。由于 \((\Psi_0)_*\partial_{\vartheta_i}=Y_i=X_{I_i}\)，有
	\(\iota_{\partial_{\vartheta_i}}\Omega=\mathrm{d}I_i\)。
	因此
	\(\Omega-\sum_i\mathrm{d}\vartheta_i\wedge\mathrm{d}I_i\)
	插入任一角方向都为零，只含 \(\mathrm{d}I_j\wedge\mathrm{d}I_k\) 项。

	每个 \(Y_i\) 是 Hamilton 场，其流保持 \(\omega\)，所以 \(\Omega\) 对全部角平移不变。刚才那些系数遂不依赖 \(\vartheta\)，得到来自 \(A\) 的二形式 \(\beta\)。两边取外微分，或直接拉回到截面 \(\vartheta=0\)，即得 \(\mathrm{d}\beta=0\)。在该截面上角微分项消失，因此 \(\beta=s^*\omega\)，这里已把 \(s\) 的参数由 \(b\) 改成 \(I\)。
\end{proof}

所以角变量的起点选择会影响剩余的横向配对。若截面本来是 Lagrange，\(\beta\) 已经为零；一般截面则需要再作一次随作用变量变化的角平移。

\subsection{辛形式的标准表达}
\label{g4:subsec:1102:04}

闭二形式 \(\beta\) 定义在开球 \(A\) 上，因而能写成一个一形式的外微分。把这个一形式的系数直接加入角变量，就能消去\eqref{g4:eq:1102:angle-base-form}中的附加项。

\begin{theorem}[Liouville--Arnold 邻域作用角定理]\label{g4:thm:1102:action-angle}
	设 \((M^{2n},\omega)\)、\(n\ge1\)，为无边界、Hausdorff 且第二可数的辛流形，\(F=(f_1,\ldots,f_n):M\to\mathbb R^n\) 光滑，且 \(\{f_i,f_j\}=0\)。设 \(L\) 是 \(F^{-1}(b_0)\) 的非空紧连通分支，\(\operatorname{rank}\mathrm{d}F=n\) 在 \(L\) 上成立。则存在 \(L\) 的开邻域 \(U\)、开球 \(A\subset\mathbb R^n\)，以及辛同胚
	\begin{equation}\label{g4:eq:1102:action-angle-map}
		\Psi:\left(\mathbb T^n\times A,\
		\sum_{i=1}^n\mathrm{d}\theta_i\wedge\mathrm{d}I_i\right)
		\longrightarrow(U,\omega),\qquad
		\mathbb T^n=(\mathbb R/2\pi\mathbb Z)^n,
	\end{equation}
	使每个 \(\mathbb T^n\times\{I\}\) 映到一个完整紧连通共同水平组件，且 \(F\Psi=b(I)\)，其中 \(b:A\to B\subset\mathbb R^n\) 为微分同胚。

	邻域上可选 \(\lambda\) 满足 \(\omega=-\mathrm{d}\lambda\)，作用变量由\eqref{g4:eq:1102:action-integral}给出，并满足\eqref{g4:eq:1102:action-derivative}。对任意 \(H=h(F)\)，方程在这些坐标中为
	\begin{equation}\label{g4:eq:1102:linear-hamilton-flow}
		\dot I_i=0,\qquad
		\dot\theta_i=\frac{\partial H}{\partial I_i}(I).
	\end{equation}
	因此每个纤维上的运动是速度恒定的环面平移。
\end{theorem}
\begin{proof}
	\Cref{g4:lem:1102:compact-fibration,g4:lem:1102:smooth-period-basis}已经构造完整环面邻域及光滑周期基底；\cref{g4:prop:1102:action-derivative}给出作用坐标。\Cref{g4:prop:1102:angle-form}把尚需处理的部分归结为 \(A\) 上的闭二形式 \(\beta\)。

	取球心 \(I^0\)，令
	\[
		\alpha_I(v)=\int_0^1t\,
		\beta_{I^0+t(I-I^0)}(I-I^0,v)\,\mathrm{d}t,
		\qquad
		\alpha=\sum_i a_i(I)\,\mathrm{d}I_i.
	\]
	这是\cref{g4:lem:0803:relative-poincare}对以点为底的向量丛所给的径向同伦公式。该点上二形式的拉回为零，所以 \(\mathrm{d}\alpha=\beta\)；这里不需要 \(\beta\) 在球心作为环境二形式为零。各 \(a_i\) 光滑。

	现在令
	\[
		\theta_i=\vartheta_i+a_i(I)\pmod{2\pi}.
	\]
	这是 \(\mathbb T^n\times A\) 的微分同胚，逆变换为
	\(\vartheta_i=\theta_i-a_i(I)\)。并且
	\[
		\sum_i\mathrm{d}\theta_i\wedge\mathrm{d}I_i
		=\sum_i\mathrm{d}\vartheta_i\wedge\mathrm{d}I_i
		 +\sum_i\mathrm{d}a_i\wedge\mathrm{d}I_i
		=\sum_i\mathrm{d}\vartheta_i\wedge\mathrm{d}I_i+\beta.
	\]
	因此 \(\Psi(\theta,I)=\Psi_0(\theta-a(I),I)\) 满足所需的辛拉回等式。角平移不改变纤维中的闭路绕行类，故前面定义的作用积分仍对应新的坐标基本圈。

	最后写
	\(X_H=\sum_i(u_i\partial_{\theta_i}+v_i\partial_{I_i})\)。
	由 \(H\) 只依赖 \(I\)，矩形坐标中的收缩方程为
	\[
		\sum_i(u_i\,\mathrm{d}I_i-v_i\,\mathrm{d}\theta_i)
		=\sum_i\frac{\partial H}{\partial I_i}\,\mathrm{d}I_i.
	\]
	所以 \(v_i=0\)、\(u_i=\partial H/\partial I_i\)，得到运动方程。其解为
	\(I(t)=I(0)\)、\(\theta(t)=\theta(0)+t\nabla_IH(I(0))\) 模 \(2\pi\)，对全部实时间有定义。
\end{proof}

这样的 \((\theta,I)\) 称为\term[zuoyongjiaozuobiao]{作用角坐标}[action-angle coordinates]。定理给出的是角变量取值于环面的邻域坐标；把每个角限制在一个普通区间内，才得到 Euclid 意义下的 Darboux 坐标图。整个环面不能被一张单值实角坐标覆盖，但圆周值角变量在这里的整个邻域上有定义。Cannas da Silva 的定理5.14陈述作用角结论，并将邻域坐标的证明指向进一步文献；上述证明则从指定紧组件出发完成了邻域、周期和角平移的构造。\cite[第5.5节，定理5.14，第152--153页]{Cannas2006Symplectic}

\begin{example}\label{g4:exm:1102:base-form-correction}
	在 \(\mathbb T^2\times A\) 上考虑
	\[
		\omega=\mathrm{d}\vartheta_1\wedge\mathrm{d}I_1
		+\mathrm{d}\vartheta_2\wedge\mathrm{d}I_2
		+\kappa\,\mathrm{d}I_1\wedge\mathrm{d}I_2,
		\qquad \kappa\in\mathbb R.
	\]
	它的 Hamilton 场 \(X_{I_i}\) 已经是 \(\partial_{\vartheta_i}\)，但截面 \(\vartheta=0\) 的拉回形式为
	\(\kappa\,\mathrm{d}I_1\wedge\mathrm{d}I_2\)。取
	\(\alpha=\kappa I_1\,\mathrm{d}I_2\)，便得到修正角
	\[
		\theta_1=\vartheta_1,\qquad
		\theta_2=\vartheta_2+\kappa I_1.
	\]
	直接求微分可见 \(\omega=\mathrm{d}\theta_1\wedge\mathrm{d}I_1+
	\mathrm{d}\theta_2\wedge\mathrm{d}I_2\)。这个例子说明，仅把场的周期归一化还不足以完成作用角坐标，截面的横向配对也必须处理。
\end{example}

\subsection{参数及格基变化}
\label{g4:subsec:1102:05}

作用变量记录的是纤维中的整数闭路与原始形式的配对；最初用哪些函数标记纤维，是另一项选择。下面分别考察重新标记参数、改选格基和改变角起点。

\begin{proposition}\label{g4:prop:1102:action-angle-changes}
	在上述作用角邻域内，有以下变换规则。
	\begin{enumerate}
		\item 若用 \(b'=h(b)\) 重新参数化，其中 \(h\) 是参数邻域之间的微分同胚，并用 \(F'=h(F)\) 的 Hamilton 场表示同一组周期闭路，则其周期矩阵为
		\begin{equation}\label{g4:eq:1102:parameter-period-change}
			\Lambda'(b')=\mathrm{D}h(b)^{-\mathsf T}\Lambda(b).
		\end{equation}
		作用积分本身不变。
		\item 在连通参数邻域内，改选光滑整数周期基底并改选满足 \(-\mathrm{d}\lambda=\omega\) 的原始形式，只会把作用变量改成
		\begin{equation}\label{g4:eq:1102:integral-affine-change}
			I'=AI+c,\qquad A\in\operatorname{GL}(n,\mathbb Z),\
			c\in\mathbb R^n,
		\end{equation}
		其中 \(A,c\) 为常量。
		\item 两组适配同一环面纤维族的作用角坐标，在缩小的连通重叠域内满足
		\begin{equation}\label{g4:eq:1102:angle-change}
			I'=AI+c,\qquad
			\theta'=A^{-\mathsf T}\theta+f(I')\pmod{2\pi}.
		\end{equation}
		对 \(f\) 的任一局部实值提升，辛条件等价于
		\[
			\mathrm{d}\left(\sum_i f_i(I')\,\mathrm{d}I'_i\right)=0.
		\]
		特别地，在足够小的开球上可写 \(f_i=\partial S/\partial I'_i\)，其中 \(S\) 为光滑函数。
	\end{enumerate}
\end{proposition}
\begin{proof}
	链式法则和 Hamilton 场的定义给出
	\(X_{f'_a}=\sum_j(\partial h_a/\partial b_j)X_{f_j}\)。
	因此新多时间 \(t'\) 对应旧多时间 \(\mathrm{D}h(b)^{\mathsf T}t'\)。同一闭路的周期列必须满足
	\(\Lambda=\mathrm{D}h^{\mathsf T}\Lambda'\)，得到第一式。并且
	\[
		\Lambda'^{\mathsf T}\mathrm{d}b'
		=\Lambda^{\mathsf T}\mathrm{D}h^{-1}\mathrm{D}h\,\mathrm{d}b
		=\Lambda^{\mathsf T}\mathrm{d}b,
	\]
	与作用积分的微分公式相容。

	若另一组周期基底为 \(\Lambda'=\Lambda C\)，则 \(C(b)\) 的每一列为整数向量，且 \(C(b)\in\operatorname{GL}(n,\mathbb Z)\)。它光滑而取值于离散集，所以在连通参数域上为常量。相应闭路的绕行类按列 \(C\) 组合；\(\lambda\) 在纤维上的限制闭，故其积分按同样整数系数组合，得到 \(I'=C^{\mathsf T}I\)。令 \(A=C^{\mathsf T}\) 即可。

	再设原始形式从 \(\lambda\) 换成 \(\lambda+\eta\)，则 \(\mathrm{d}\eta=0\)。沿参数变化的闭路族作与\cref{g4:prop:1102:action-derivative}相同的积分求导，右端变成 \(\int\mathrm{d}\eta(V,T)=0\)。因此各个周期积分的改变量都是常数，给出 \(c\)。改变闭路的起点只产生同一纤维中的同伦，对积分没有影响。

	两组作用角坐标各自的基本圈给出完整整数基底，各自的典范原始形式 \(\sum_i I_i\,\mathrm{d}\theta_i\)、\(\sum_i I'_i\,\mathrm{d}\theta'_i\) 的基本圈积分恰为 \(2\pi I_i\)、\(2\pi I'_i\)，故第二项给出 \(I'=AI+c\)。由
	\[
		X_{I'_i}=\sum_j A_{ij}X_{I_j},\qquad
		X_{I_j}=\partial_{\theta_j},\qquad
		X_{I'_i}=\partial_{\theta'_i},
	\]
	得到 \(\partial\theta'/\partial\theta=A^{-\mathsf T}\)。所以
	\(\theta'-A^{-\mathsf T}\theta\) 沿每个连通纤维恒定，是参数的圆周值函数 \(f(I')\)。局部选其实值提升后，代入二形式等式有
	\[
		\sum_i\mathrm{d}\theta'_i\wedge\mathrm{d}I'_i
		=\sum_i\mathrm{d}\theta_i\wedge\mathrm{d}I_i
		 +\sum_i\mathrm{d}f_i\wedge\mathrm{d}I'_i.
	\]
	因此辛条件就是所列闭性。不同提升相差 \(2\pi\) 的整数常向量，不改变这个条件。在小开球上，闭一形式由径向同伦得到原始函数 \(S\)，遂有 \(f_i=\partial S/\partial I'_i\)。
\end{proof}

格基变化只能使用整数可逆矩阵。任意实线性变换虽然能重新组合微分独立的第一积分，却未必把基本闭路变成整数基底；若只取某个真子格，得到的角参数会覆盖同一环面多次。作用角坐标同时保留了辛配对和闭路的整周期信息。下一节将考察这些局部格基沿较大参数空间移动时，为什么可能无法整体相容。
% 纲要条目排为编号小节标题，并进入目录与书签；说明段仍保留于纲要注释中。

% !TeX root = ../../Book4.tex
% 纲要标题：### 11.3 整体障碍
% 纲要位置：几何学计划纲要.md，第 3900 行。
% 正文状态：初稿；尚未正式审核。

% 纲要内容开始
% 
% * 周期格局部系统
% * 整仿射结构
% * 单值矩阵和不可整体选取作用角坐标
% * 焦点型奇纤维的代表模型
% 
% 纲要内容结束

\section{整体障碍}
\label{g4:sec:1103}

一个正则环面的附近已经有作用角坐标。若沿底空间移动这个环面，每经过一个小邻域都可以继续选择坐标；回到出发点时，原来的基本闭路却可能变成另一组整数基。本节先把这种变化记录为周期格的单值，再用一个显式粘合构造算出非平凡的单值矩阵。

\subsection{周期格局部系统}
\label{g4:subsec:1103:01}
设 \(\pi:(M^{2n},\omega)\to B^n\) 是光滑满秩映射，每条纤维都是紧连通 Lagrange 子流形，且局部具有前节构造的作用角邻域。这样的映射称为\term[lagrangehuanmianxianweihua]{Lagrange 环面纤维化}[Lagrangian torus fibration]。对完全可积系统，可以在正则值附近只选定一族连通纤维；这里允许把这些局部族沿一般底流形粘合起来。

\begin{definition}\label{g4:def:1103:period-covectors}
	对 \(b\in B\) 和 \(\alpha\in\mathrm{T}_b^*B\)，选局部函数 \(h\) 满足 \((\mathrm{d}h)_b=\alpha\)。向量场 \(X_{h\circ\pi}\) 在 \(L_b=\pi^{-1}(b)\) 上只由 \(\alpha\) 决定，记为 \(Y_\alpha\)。定义余切空间中的周期格\termidx[zhouqige]{周期格}
	\[
		\mathcal P_b=\{\alpha\in\mathrm{T}_b^*B:\text{\(Y_\alpha\) 的时间 \(2\pi\) 流在 \(L_b\) 上为恒等}\}.
	\]
\end{definition}

这个定义把周期统一归一化为 \(2\pi\)。在作用角坐标 \((\theta,I)\) 中，若 \(\alpha=\sum_j a_j\,\mathrm{d}I_j\)，则 \(Y_\alpha=\sum_j a_j\partial_{\theta_j}\)，所以
\begin{equation}\label{g4:eq:1103:period-frame}
	\mathcal P_b=\bigoplus_{j=1}^n\mathbb Z(\mathrm{d}I_j)_b.
\end{equation}
纤维上的这些场可交换且完备。沿 \(Y_\alpha\) 走时间 \(2\pi\) 得到闭路，给出 \(\mathcal P_b\cong H_1(L_b;\mathbb Z)\)：在角坐标中，这就是把整数向量送到相应的绕行次数。改变闭路起点只作环面平移；平移沿一条道路同伦于恒等，故同调类不变。

\begin{proposition}\label{g4:prop:1103:local-system}
	局部周期格基之间的过渡矩阵属于 \(\operatorname{GL}(n,\mathbb Z)\)，并在每个连通的重叠区域上为常矩阵。因此周期格可以沿道路延拓；固定起点与终点时，延拓只依赖道路的端点固定同伦类。
\end{proposition}
\begin{proof}
	两组光滑周期格基在每一点都是同一个自由整数模的基，所以换基矩阵及其逆的系数都是整数。系数同时随底点连续变化，因而在连通重叠区域上恒定。三组基之间的换基满足通常的矩阵余循环等式。

	把一条道路分成有限段，使每段落在一个局部坐标域内，并在这些域内保持整数坐标不变，就得到了延拓。细分不改变结果，更换覆盖时用共同细分及换基等式比较。对一个端点固定的道路同伦，把参数方形剖分成足够小的矩形，使每个矩形的像落在一个坐标域内；这样的剖分由紧性和有限开覆盖得到。每个小矩形边界的延拓为恒等，相邻边相消后得到同伦两边的延拓相同。
\end{proof}

像这样由局部固定的整数模及局部常值过渡矩阵组成的对象，称为整数\term[jubuxitong]{局部系统}[local system]。这里的“局部常值”是指格中的整数标号可以局部固定，余切向量本身仍随底点光滑变化。

\subsection{整仿射结构}
\label{g4:subsec:1103:02}
\begin{proposition}\label{g4:prop:1103:integral-affine}
	环面纤维化的局部作用变量在连通重叠区域上满足
	\begin{equation}\label{g4:eq:1103:affine-transition}
		I'=AI+c,\qquad A\in\operatorname{GL}(n,\mathbb Z),\quad c\in\mathbb R^n,
	\end{equation}
	其中 \(A,c\) 都是常量。因此底空间自然具有\term[zhengfangshejiegou]{整仿射结构}[integral affine structure]，即坐标变换为整数可逆线性变换后接平移的坐标图册。
\end{proposition}
\begin{proof}
	由\cref{g4:prop:1103:local-system}，两组周期格基满足 \(\mathrm{d}I'=A\,\mathrm{d}I\)，其中 \(A\) 为常整数可逆矩阵。故 \(\mathrm{d}(I'-AI)=0\)，连通性给出常向量 \(c\)。各 \(I\) 确为局部坐标，已由作用角定理中的微分独立性保证。
\end{proof}

整数矩阵来自基本闭路的换基，平移常数来自作用积分原点的选择。如果在同一邻域内把原始一形式 \(\lambda\) 换成 \(\lambda+\eta\)，其中 \(\mathrm{d}\eta=0\)，同一族基本闭路上的 \(\int\eta\) 不随底点改变：两条相邻闭路之间的圆柱用 Stokes 公式给出积分差为零。这正是作用变量允许加常数的原因。

\begin{example}\label{g4:exm:1103:affine-circle}
	在二维环面上令 \(\theta,s\in\mathbb R/(2\pi\mathbb Z)\)，取
	\[
		\omega=\mathrm{d}\theta\wedge\mathrm{d}s,\qquad \pi(\theta,s)=s.
	\]
	这里 \(\mathrm{d}s\) 表示圆上局部实提升的微分；不同提升相差 \(2\pi\) 的整数倍，故微分全局一致。纤维的周期格有全局基 \(\mathrm{d}s\)，线性换基完全平凡。不过不存在全局实值作用变量 \(I\) 使 \(\mathrm{d}I\) 成为这个格的基。否则在连通底圆上必有 \(\mathrm{d}I=\mathrm{d}s\) 或 \(-\mathrm{d}s\)，沿底圆积分分别为 \(2\pi\) 或 \(-2\pi\)，与恰当形式的闭路积分为零矛盾。
\end{example}

这个例子区分了两件事：周期格能否选取全局基，以及基中的闭一形式能否同时积分成全局实值函数。即使这两步都能完成，选取全局角变量还须使各纤维中的角原点相容。下面的例子中，作用变量已经全局存在，但底空间上的闭二形式阻止了这种角原点选择。

\begin{example}\label{g4:exm:1103:magnetic-obstruction}
	令 \(B=\mathbb R^3\setminus\{0\}\)，底坐标为 \(b=(b_1,b_2,b_3)\)，取光滑二形式
	\[
		\beta=\frac{b_1\,\mathrm{d}b_2\wedge\mathrm{d}b_3+
		b_2\,\mathrm{d}b_3\wedge\mathrm{d}b_1+
		b_3\,\mathrm{d}b_1\wedge\mathrm{d}b_2}{|b|^3}.
	\]
	记 \(\mathbb T^3=(\mathbb R/2\pi\mathbb Z)^3\)。在 \(M=\mathbb T^3\times B\) 上定义
	\[
		\omega=\sum_{j=1}^3\mathrm{d}\theta_j\wedge\mathrm{d}b_j+\pi^*\beta,
		\qquad \pi(\theta,b)=b.
	\]
	其中 \(\mathrm{d}\theta_j\) 表示第 \(j\) 个角圆上标准闭一形式的拉回，并不要求选择角的全局实提升。

	先核对这个模型的辛结构与周期格。在 \(b\ne0\) 处，
	\[
		\operatorname{div}\frac b{|b|^3}
		=\frac3{|b|^3}-\frac{3|b|^2}{|b|^5}=0,
	\]
	故 \(\mathrm{d}\beta=0\)，从而 \(\mathrm{d}\omega=0\)。若切向量 \((u,v)\) 与所有纯角向量的辛配对为零，则 \(v=0\)；再与所有纯底向量配对，得到 \(u=0\)。所以 \(\omega\) 非退化。又有
	\[
		\iota_{\partial_{\theta_j}}\omega=\mathrm{d}b_j,
		\qquad X_{b_j}=\partial_{\theta_j}.
	\]
	因此 \(b_1,b_2,b_3\) 处处独立且两两 Poisson 对合。每条纤维是 Lagrange 三维环面，角场的流为平移且周期为 \(2\pi\)。按本节的时间 \(2\pi\) 归一化，\(\sum_j a_j(\mathrm{d}b_j)_b\) 属于周期格当且仅当全部 \(a_j\) 都是整数，故
	\[
		\mathcal P_b=\bigoplus_{j=1}^3\mathbb Z(\mathrm{d}b_j)_b.
	\]
	这个格已有全局恰当基，\(b_j\) 本身就是相应的全局实值作用变量。

	现在检验角原点能否相容。任一光滑截面写作 \(s(b)=(f(b),b)\)，其中 \(f:B\to\mathbb T^3\)。记 \(\alpha_j=f^*(\mathrm{d}\theta_j)\)；各 \(\alpha_j\) 都是全局闭一形式，无须提升 \(f\) 的角分量。于是
	\[
		s^*\omega=\sum_j\alpha_j\wedge\mathrm{d}b_j+\beta
		=\beta-\mathrm{d}\!\left(\sum_jb_j\alpha_j\right).
	\]
	在按外法向定向的单位球面 \(S^2\subset B\) 上，\(\beta\) 就是面积形式，因而积分为 \(4\pi\)。由闭曲面上的 Stokes 定理，恰当形式的积分为零，所以
	\[
		\int_{S^2}s^*\omega=4\pi.
	\]
	故不存在满足 \(s^*\omega=0\) 的截面，即 \(\pi\) 没有 Lagrange 截面。若存在与 \(\pi\) 相容的全局作用角坐标，把全部角固定为零就会产生一个 Lagrange 截面，与此矛盾。非恰当二形式 \(\beta\) 由此给出了周期格以外的整体障碍。
\end{example}

\subsection{单值矩阵和不可整体选取作用角坐标}
\label{g4:subsec:1103:03}
固定底点 \(b_0\) 和纤维的一组基本同调类 \(\gamma_1,\ldots,\gamma_n\)。沿以 \(b_0\) 为端点的环路 \(\ell\) 延拓后，记所得基为 \(\widetilde\gamma_1,\ldots,\widetilde\gamma_n\)，用约定
\begin{equation}\label{g4:eq:1103:monodromy-matrix}
	(\widetilde\gamma_1,\ldots,\widetilde\gamma_n)
	=(\gamma_1,\ldots,\gamma_n)C_\ell,
	\qquad C_\ell\in\operatorname{GL}(n,\mathbb Z)
\end{equation}
定义\term[hamiltondanzhi]{Hamilton 单值}[Hamiltonian monodromy]的矩阵。这里的单值就是沿环路延拓周期格所产生的整数换基。若先沿 \(\ell\) 再沿 \(m\)，延拓算子的复合为 \(T_m\circ T_\ell\)，在上述列坐标约定下 \(C_{\ell*m}=C_mC_\ell\)。这与第一卷“先走左侧道路”的拼接约定一致；采用相反的拼接顺序即得到通常写作的左表示。

\begin{proposition}\label{g4:prop:1103:monodromy-obstruction}
	设底空间 \(B\) 连通。周期格局部系统有一组全局局部常值基，当且仅当所有环路的单值矩阵都是单位矩阵。特别地，只要某个 \(C_\ell\ne I_n\)，就不能在整个纤维化上选取全局作用角坐标。
\end{proposition}
\begin{proof}
	全局基沿任何道路延拓后仍是同一组基，所以单值必平凡。反过来，若所有环路的单值均平凡，从 \(b_0\) 选道路到任意点并延拓初始基；两条道路的延拓之差由一个基点环路给出，因而结果无关道路。在每个局部平凡化中，所得基保持同一组整数坐标，故是所需的全局基。

	全局作用角坐标的各个角圆给出一组全局基本闭路，等价地，全局作用变量的微分给出周期格基。因此非平凡单值排除这种坐标。
\end{proof}

改变初始基把全部矩阵同时共轭，所以“是否为单位矩阵”与换基无关。反向绕行把矩阵变成逆矩阵。若闭路基按 \(C\) 变化，则由作用积分定义，相应作用列向量按 \(C^{\mathsf T}\) 变化，另可有常向量平移；这两种矩阵不可混写。例如
\[
	C=\begin{pmatrix}1&1\\0&1\end{pmatrix}
	\quad\Longleftrightarrow\quad
	\widetilde\gamma_1=\gamma_1,\quad
	\widetilde\gamma_2=\gamma_2+\gamma_1.
\]
一个基本圆保持不变，另一个多绕这个圆一次。下一小节直接构造具有这一现象的纤维化。

\subsection{焦点型奇纤维的代表模型}
\label{g4:subsec:1103:04}
先在 \(\mathbb R^4\) 上取坐标 \((x_1,x_2,y_1,y_2)\)、辛形式 \(\omega_0=\sum_j\mathrm{d}x_j\wedge\mathrm{d}y_j\)，并令
\begin{equation}\label{g4:eq:1103:focus-quadratic}
	a=Q_1=x_1y_1+x_2y_2,\qquad
	b=Q_2=x_1y_2-x_2y_1.
\end{equation}
这对二次函数称为标准\term[jiaodianxingqidian]{焦点型奇点}[focus-focus singularity]模型。其 Hamilton 场分别为
\[
	X_{Q_1}=(x,-y),\qquad
	X_{Q_2}=(-x_2,x_1,-y_2,y_1).
\]
前者使 \(x\) 膨胀、\(y\) 收缩，后者使两个平面同时正向旋转。旋转保持内积 \(x\cdot y\)，故 \(\{Q_1,Q_2\}=X_{Q_2}Q_1=0\)。对实数 \(u,v\)，\(X_{uQ_1+vQ_2}\) 的复特征值为 \(u\pm\mathrm{i}v,-u\pm\mathrm{i}v\)，显示出伸缩与旋转同时出现的线性行为。

\ProseParagraphBreak

写 \(z=x_1+\mathrm{i}x_2\)、\(w=y_1+\mathrm{i}y_2\)，则
\[
	c=a+\mathrm{i}b=\overline z w.
\]
若 \((z,w)\ne(0,0)\)，微分 \(\mathrm{d}c=w\,\mathrm{d}\overline z+\overline z\,\mathrm{d}w\) 满射到 \(\mathbb C\)，所以唯一临界点为原点。在 \(z\ne0\) 的区域令 \(t=\log|z|\)、\(\vartheta=\arg z\)。由 \(w=c/\overline z\) 直接计算得
\begin{equation}\label{g4:eq:1103:focus-cylinder}
	\lambda_0=\sum_jy_j\,\mathrm{d}x_j
	=a\,\mathrm{d}t+b\,\mathrm{d}\vartheta,
	\qquad
	\omega_0=-\mathrm{d}\lambda_0
	=\mathrm{d}t\wedge\mathrm{d}a+\mathrm{d}\vartheta\wedge\mathrm{d}b.
\end{equation}
因此 \((Q_1,Q_2)\) 的共同流就是 \((t,\vartheta)\) 的平移。单独的二次模型具有非紧正则纤维；为看见环面上的单值，还要把圆柱的两端辛地接起来。

\begin{proposition}\label{g4:prop:1103:focus-gluing}
	取 \(0<\delta<1\)。有一个辛四维流形 \(M_\delta\) 和固有可积映射 \(F:M_\delta\to B_\delta=\{c\in\mathbb C:|c|<\delta\}\)，其非零纤维为连通二维环面，零纤维只有一个临界点且在该点附近等于模型\eqref{g4:eq:1103:focus-quadratic}。
\end{proposition}
\begin{proof}
	在 \(\mathbb C^2\) 中取
	\[
		P_\delta=\{(z,w):|z|\le1,\ |w|\le1,\ |\overline z w|<\delta\}.
	\]
	它有两个互不相交的边界面 \(A=\{|z|=1\}\cap P_\delta\) 和 \(D=\{|w|=1\}\cap P_\delta\)，因为同时有 \(|z|=|w|=1\) 将违反 \(|c|<\delta\)。用
	\begin{equation}\label{g4:eq:1103:focus-gluing-map}
		\Psi(z,w)=\left(z^2\overline w,\frac1{\overline z}\right)
	\end{equation}
	把 \(A\) 与 \(D\) 识别。这个公式在 \(A\) 的邻域光滑，其逆在 \(D\) 的邻域为
	\[
		(z',w')\longmapsto
		\left(\frac1{\overline{w'}},\overline{z'}(w')^2\right).
	\]
	它保持 \(\overline z w\)，并将 \(A\) 内侧 \(|z|<1\) 映到 \(D\) 外侧 \(|w'|>1\)。因而两侧的领圈通过此公式拼成跨越接缝的光滑坐标图。

	核验辛性时先设 \(c\ne0\)，写 \(r=|c|\)、\(\varphi=\arg c\)。由公式得
	\[
		t'=t+\log r,\qquad \vartheta'=\vartheta-\varphi.
	\]
	注意
	\[
		a\,\mathrm{d}\log r-b\,\mathrm{d}\varphi=\mathrm{d}a,
	\]
	所以 \(\Psi^*\lambda_0=\lambda_0+\mathrm{d}a\)，从而 \(\Psi^*\omega_0=\omega_0\)。两边在领圈内都光滑，该等式由稠密性延伸到 \(c=0\)。辛形式因此在商 \(M_\delta=P_\delta/{\sim}\) 上粘合。

	此商也是 Hausdorff 且第二可数的流形。具体说，边界两面在 \(P_\delta\) 中闭，任意闭集的饱和是该集与两面的对应像之并，仍为闭集，所以商映射是闭映射。每个等价类至多含两点；在度量空间 \(P_\delta\) 中分开两个有限等价类，再用闭映射缩成饱和邻域，可分开任意两个商点。内部图与跨接缝的领圈图覆盖商，并可从两面和内部的可数坐标覆盖选出可数图册。这也验证了商所需的分离性与可数性。

	函数 \(c=\overline z w\) 下降为 \(F\)。紧集 \(K\subset B_\delta\) 的原像是 \(\{|z|,|w|\le1,\ \overline z w\in K\}\) 的商；括号内是 \(\mathbb R^4\) 中的闭有界集，故原像紧。唯一临界点仍是内部的原点。对 \(c\ne0\)，纤维在识别前为
	\[
		\log|c|\le t\le0,\qquad \vartheta\in\mathbb R/(2\pi\mathbb Z),
	\]
	其两边按常旋转 \(\vartheta\mapsto\vartheta-\arg c\) 识别。令 \(\tau=-\log|c|>0\)，在圆柱上把角改成 \(\vartheta-(\arg c)t/\tau\) 即消去边界扭转，因此商为二维环面。原点邻域没有参与粘合，仍具有原二次模型。
\end{proof}

零纤维在粘合前是两个闭单位圆盘 \(\{w=0\}\) 与 \(\{z=0\}\)，圆心为同一点；边界识别为 \((z,0)\sim(0,z)\)。若先把两圆盘边界粘合，就得到球面，再将两个圆心识别。所得空间也可以由环面把一条基本圆压成一点得到，称为\term[zhensuohuanmian]{针缩环面}[pinched torus]。因此这里的奇纤维确有可见的拓扑退化。

\begin{proposition}\label{g4:prop:1103:focus-monodromy}
	上述模型在穿孔圆盘 \(0<|c|<\delta\) 上的周期格具有非平凡单值。按正向绕行及\eqref{g4:eq:1103:monodromy-matrix}的约定，合适的基本闭路基给出
	\[
		C=\begin{pmatrix}1&1\\0&1\end{pmatrix}.
	\]
\end{proposition}
\begin{proof}
	选定 \(\varphi=\arg c\) 的一个局部分支。共同流的时间参数为 \((t,\vartheta)\)，一个周期是 \(\ell_1=(0,2\pi)\)。边界识别的平移向量为 \((\log|c|,-\varphi)\)，所以另一周期可取
	\begin{equation}\label{g4:eq:1103:focus-periods}
		\ell_2=(-\log|c|,\varphi).
	\end{equation}
	这两者生成全部时间周期：任意闭合的平移跨越整数次接缝，减去相应倍数的 \(\ell_2\) 后，其 \(t\) 位移为零，只能再绕整数次角圆。二者线性独立，因为 \(-\log|c|>0\)。

	沿底圆正向一周，\(\varphi\) 增加 \(2\pi\)。于是 \(\ell_1\) 不变而 \(\ell_2\) 变成 \(\ell_2+\ell_1\)，对应的基本闭路恰按所述矩阵变化。这个矩阵不是单位矩阵，因而\cref{g4:prop:1103:monodromy-obstruction}适用。
\end{proof}

周期还可以直接积分成局部作用变量。按前节的 \(2\pi\) 归一化，取
\begin{equation}\label{g4:eq:1103:focus-actions}
	I_1=b,\qquad
	I_2=\frac{-a\log|c|+b\arg c+a}{2\pi}.
\end{equation}
由上面的一形式恒等式，
\[
	\mathrm{d}I_2=\frac{-\log|c|\,\mathrm{d}a+\arg c\,\mathrm{d}b}{2\pi}.
\]
这与时间周期 \(\ell_2\) 完全对应，且 \(\mathrm{d}I_1,\mathrm{d}I_2\) 独立。绕行后 \(I_1'=I_1\)、\(I_2'=I_2+I_1\)，清楚地表现出作用变量的多值性。图\ref{g4:fig:1103:focus-monodromy}记录这个换基。

% !TeX root = ../Book4.tex
\begin{figure}[htbp]
	\centering
	\begin{tikzpicture}[font=\small,>=stealth,scale=0.95]
		\draw[->] (-1.8,0)--(1.8,0) node[right] {\(a\)};
		\draw[->] (0,-1.5)--(0,1.7) node[above] {\(b\)};
		\draw[thick,->] (25:1.1) arc[start angle=25,end angle=345,radius=1.1];
		\fill (0,0) circle (2pt);
		\node[below left] at (0,0) {\(0\)};
		\node[align=center] at (0,-2) {绕奇值正向一周};
		\begin{scope}[shift={(4.5,0)}]
			\draw[step=0.65,gray!30] (-0.7,-0.7) grid (2,1.4);
			\draw[->,very thick] (0,0)--(1.3,0) node[below] {\(\gamma_1\)};
			\draw[->,very thick] (0,0)--(0,1.3) node[left] {\(\gamma_2\)};
			\draw[->,very thick,dashed] (0,0)--(1.3,1.3) node[right] {\(\widetilde\gamma_2\)};
			\node[align=center] at (0.7,-2) {\(\widetilde\gamma_1=\gamma_1\)\\\(\widetilde\gamma_2=\gamma_2+\gamma_1\)};
		\end{scope}
	\end{tikzpicture}
	\caption[焦点型模型的周期换基]{左图是底空间的绕行道路，右图在同一个整数同调格中比较绕行前后的基；虚线表示延拓后的第二个基本闭路。右图画的是同调类，并非纤维中的实际轨道形状。}
	\label{g4:fig:1103:focus-monodromy}
\end{figure}


本节的构造同时给出了辛粘合、紧纤维与周期计算，因而无需借助一般奇点分类即可检验单值障碍。对一个给定系统，若只算出某平衡点的二次部分等于\eqref{g4:eq:1103:focus-quadratic}，还须研究该系统完整的奇纤维邻域，才能把这里的单值结论转移过去。下一节的球摆计算将分别展示二次部分和实际奇纤维的形状。
% 纲要条目排为编号小节标题，并进入目录与书签；说明段仍保留于纲要注释中。

% !TeX root = ../../Book4.tex
% 纲要标题：### 11.4 典型可积系统
% 纲要位置：几何学计划纲要.md，第 3907 行。
% 正文状态：初稿；尚未正式审核。

% 纲要内容开始
% 
% * 振子与环面平移
% * 测地流的对称性
% * 球摆及奇异纤维的入口
% 
% ---
% 
% 纲要内容结束

\section{典型可积系统}
\label{g4:sec:1104}

作用角坐标把正则紧纤维附近的运动化为环面平移。具体系统仍有三件事需要计算：哪些函数构成独立的交换积分，哪些共同水平集满足紧性与正则性，以及频率怎样由作用变量决定。振子可以把坐标与频率完全写出；测地流显示对称性怎样提供积分；球摆则让正则环面与奇异纤维出现在同一个能量图景中。本节沿用角变量周期为 \(2\pi\)、\(\omega=\sum_j\mathrm{d}\theta_j\wedge\mathrm{d}I_j\) 的约定。

\subsection{振子与环面平移}
\label{g4:subsec:1104:01}
\cref{g4:exm:0902:two-oscillators}已经证明，两个振子的能量分量彼此对合，并在两分量都非零处独立。现在计算它们的作用积分，并辨认环面上的轨道。

\begin{example}\label{g4:exm:1104:oscillator-actions}
	在 \(\mathbb R^{2n}\) 上取
	\[
		\omega=\sum_{j=1}^n\mathrm{d}q_j\wedge\mathrm{d}p_j,
		\qquad I_j=\frac{q_j^2+p_j^2}{2},
		\qquad H=\sum_{j=1}^n\nu_jI_j,
		\quad \nu_j>0.
	\]
	在每个 \(I_j>0\) 的区域定义
	\begin{equation}\label{g4:eq:1104:oscillator-angles}
		q_j=\sqrt{2I_j}\cos\theta_j,
		\qquad p_j=-\sqrt{2I_j}\sin\theta_j,
		\qquad \theta_j\in\mathbb R/(2\pi\mathbb Z).
	\end{equation}
	对单个分量求微分，得到
	\[
		p_j\,\mathrm{d}q_j
		=2I_j\sin^2\theta_j\,\mathrm{d}\theta_j
		-\sin\theta_j\cos\theta_j\,\mathrm{d}I_j
		=I_j\,\mathrm{d}\theta_j
		-\mathrm{d}(I_j\sin\theta_j\cos\theta_j).
	\]
	因此典范一形式 \(\lambda=\sum_jp_j\,\mathrm{d}q_j\) 满足
	\begin{equation}\label{g4:eq:1104:oscillator-action-integral}
		\omega=-\mathrm{d}\lambda
		=\sum_j\mathrm{d}\theta_j\wedge\mathrm{d}I_j,
		\qquad
		\frac1{2\pi}\int_{\gamma_j}\lambda=I_j,
	\end{equation}
	其中 \(\gamma_j\) 固定所有作用及其余角变量，只让 \(\theta_j\) 增加 \(2\pi\)。这核对了角变量中负号的作用：沿实际运动，\(\theta_j\) 递增，作用积分为正。

	Hamilton 方程为 \(\dot I_j=0\)、\(\dot\theta_j=\nu_j\)。故能量分量 \(E_j=\nu_jI_j\)，对应频率为 \(\partial H/\partial I_j=\nu_j\)，周期为 \(2\pi/\nu_j\)。若从质量为 \(m\)、频率为 \(\nu\) 的物理振子
	\[
		H(q,p)=\frac{p^2}{2m}+\frac{m\nu^2q^2}{2}
	\]
	出发，先作辛缩放 \(Q=\sqrt{m\nu}\,q\)、\(P=p/\sqrt{m\nu}\)，便有 \(H=\nu(Q^2+P^2)/2\)。所以它的作用为 \(I=E/\nu\)，同样等于能量椭圆内部面积除以 \(2\pi\)。
\end{example}

更一般地，在作用角区域中取光滑函数 \(H=h(I)\)。收缩方程直接给出
\begin{equation}\label{g4:eq:1104:torus-translation}
	I(t)=I(0),\qquad
	\theta(t)=\theta(0)+t\,\nu(I(0)),
	\qquad \nu(I)=\nabla h(I).
\end{equation}
给定环面上的频率向量 \(\nu\ne0\)，轨道闭合当且仅当存在 \(T>0\) 使 \(T\nu\in2\pi\mathbb Z^n\)。必要性来自返回时各角的增量，充分性则由显式解得到。若存在非零整数向量 \(k\) 使 \(k\cdot\nu=0\)，称这些频率满足一个\term[gongzhen]{共振}[resonance]关系；此时 \(k\cdot\theta\) 模 \(2\pi\) 沿运动不变。

\begin{proposition}\label{g4:prop:1104:two-torus-orbits}
	设 \(\nu_1\nu_2\ne0\)。在二维环面 \((\mathbb R/2\pi\mathbb Z)^2\) 上，平移流 \(\theta(t)=\theta(0)+t(\nu_1,\nu_2)\) 的轨道在 \(\nu_2/\nu_1\in\mathbb Q\) 时均闭合，在 \(\nu_2/\nu_1\notin\mathbb Q\) 时均稠密。
\end{proposition}
\begin{proof}
	闭合条件等价于 \(T\nu_j=2\pi m_j\)，其中 \(m_j\) 为整数。因此闭合蕴含有理比；反过来，若 \(\nu_2/\nu_1=m_2/m_1\)，可以选取正的 \(T\) 使两式同时成立。

	设 \(\alpha\notin\mathbb Q\)。先证明 \(\{n\alpha\bmod1:n\in\mathbb Z\}\) 在圆上稠密。将 \(0,\alpha,\ldots,N\alpha\) 的小数部分放入 \(N\) 个长度为 \(1/N\) 的区间。两个落在同一区间，故存在非零整数 \(m\) 及整数 \(\ell\)，使 \(0<|m\alpha-\ell|<1/N\)。把符号调整后，记此正数为 \(\delta\)。圆上从 \(0\) 到 \(\lfloor1/\delta\rfloor\delta\) 的这些点都是整数倍 \(\alpha\) 的剩余类，连续间隙及最后回到 \(0\) 的间隙都不大于 \(\delta\)。令 \(N\) 任意增大，即得稠密性。

	现在指定目标角 \((a,b)\)。在时刻
	\[
		t_n=\frac{a-\theta_1(0)+2\pi n}{\nu_1},\qquad n\in\mathbb Z,
	\]
	第一角恰为 \(a\)，第二角模 \(2\pi\) 等于一个固定常数加 \(2\pi n\nu_2/\nu_1\)。上述圆上稠密性说明它可以任意接近 \(b\)。因此轨道与每个非空开矩形相交，从而在整个环面中稠密。
\end{proof}

二维的一个非零整数共振关系已经迫使频率比有理；在更高维，一个共振只把轨道限制在较低维的子环面中。例如频率 \((1,\sqrt2,1)\) 满足第一与第三分量相等，轨道仍因前两分量而不闭合。频率的算术性质与共同水平集的拓扑也应区分：同一个光滑环面既可以承载闭轨道流，也可以承载稠密轨道流。

\subsection{测地流的对称性}
\label{g4:subsec:1104:02}
设 \((Q,g)\) 为无边界光滑黎曼流形。在 \(\mathrm{T}^*Q\) 上取典范辛形式，并定义动能
\[
	H(x,p)=\frac12g^{ij}(x)p_ip_j.
\]
这个 Hamilton 系统在底空间上的投影就是带仿射参数的测地线。为将两种描述接起来，写出方程
\[
	\dot x^i=g^{ij}p_j,
	\qquad \dot p_i=-\frac12\partial_i g^{jk}\,p_jp_k
	=\frac12\partial_i g_{ab}\,\dot x^a\dot x^b.
\]
最后一个等式由对 \(g^{ja}g_{ak}=\delta^j_k\) 求导得到。再对 \(p_i=g_{ij}\dot x^j\) 求时间导数，便有
\[
	g_{ij}\ddot x^j+\partial_k g_{ij}\,\dot x^k\dot x^j
	-\frac12\partial_i g_{ab}\,\dot x^a\dot x^b=0.
\]
将中间项按 \(k,j\) 对称化并乘以逆度量，得到
\begin{equation}\label{g4:eq:1104:geodesic-hamilton-equivalence}
	\ddot x^\ell+
	\frac12g^{\ell i}
	\bigl(\partial_jg_{ki}+\partial_kg_{ji}-\partial_ig_{jk}\bigr)
	\dot x^j\dot x^k=0.
\end{equation}
括号中的系数正是第三卷第1.2节式(1.18)的度量联络系数，故这就是第三卷定义2.1.1中的测地线方程。反过来，由满足此方程的曲线取 \(p_i=g_{ij}\dot x^j\)，上述计算倒推给出 Hamilton 方程。因而两种动力学经度量同构对应；余切丛上的这种流也称为\term[cediliu]{测地流}[geodesic flow]。
% 跨卷实读：g3:thm:levi-civita（定理1.2.5）及其Koszul推导、构造证明；g3:eq:christoffel-metric（式1.18）；g3:def:geodesic（定义2.1.1）和g3:eq:geodesic-equation（式2.2）。

\begin{example}\label{g4:exm:1104:flat-torus}
	在平坦环面 \(Q=(\mathbb R/2\pi\mathbb Z)^2\)、\(g=\mathrm{d}x_1^2+\mathrm{d}x_2^2\) 上，余切丛全局写为 \(Q\times\mathbb R^2\)，且
	\[
		H=\frac{p_1^2+p_2^2}{2},\qquad
		\dot x_j=p_j,\qquad \dot p_j=0.
	\]
	平移对称性提供两个处处独立的交换积分 \(p_1,p_2\)。共同水平集 \(p=c\) 是紧环面，\(x_j\) 本身就是角变量，作用为 \(I_j=p_j\)，频率为 \(c_j\)。当 \(c_1c_2\ne0\) 时，前一命题给出测地线闭合或稠密的准确判据；一个分量为零而另一个非零时轨道闭合，两者都为零时轨道为一点。这里 \(H\) 是交换积分的函数；若改用 \((H,p_2)\)，其微分只在 \(p_1\ne0\) 处独立。
\end{example}

\begin{example}\label{g4:exm:1104:surface-revolution}
	旋转曲面在避开旋转轴的区域可以写成
	\[
		(s,\theta)\longmapsto
		\bigl(r(s)\cos\theta,r(s)\sin\theta,z(s)\bigr),
		\qquad r(s)>0,\quad r'(s)^2+z'(s)^2=1,
	\]
	其中 \(s\) 是母线的弧长。诱导度量为 \(\mathrm{d}s^2+r(s)^2\mathrm{d}\theta^2\)。在余切坐标 \((s,\theta,p_s,j)\) 中，动能及旋转的矩为
	\begin{equation}\label{g4:eq:1104:revolution-hamiltonian}
		H=\frac12\left(p_s^2+\frac{j^2}{r(s)^2}\right),
		\qquad J=j,
		\qquad \{H,J\}=\partial_\theta H=0.
	\end{equation}
	方程具体为
	\[
		\dot s=p_s,\qquad
		\dot\theta=\frac{j}{r(s)^2},\qquad
		\dot p_s=\frac{j^2r'(s)}{r(s)^3},\qquad
		\dot j=0.
	\]
	由于
	\[
		\mathrm{d}H
		=p_s\,\mathrm{d}p_s-\frac{j^2r'(s)}{r(s)^3}\,\mathrm{d}s
		+\frac{j}{r(s)^2}\,\mathrm{d}j,
	\]
	在 \(H>0\) 上，\(\mathrm{d}H,\mathrm{d}J\) 线性相关当且仅当 \(p_s=0\)、\(r'(s)=0\)。特别地，它们在稠密开集 \(p_s\ne0\) 上独立。这证明了该旋转区域中的完全可积性，并具体指出了独立性可能失效的位置。

	对单位速测地线，令 \(\beta\) 表示速度与正向母线的夹角，正向纬线的单位切向为 \(r^{-1}\partial_\theta\)。于是 \(\sin\beta=r\dot\theta\)，角动量守恒写成
	\begin{equation}\label{g4:eq:1104:revolution-conservation}
		r(s)\sin\beta=j.
	\end{equation}
	这称为\term[clairautguanxi]{Clairaut 关系}[Clairaut's relation]。\footnote{Alexis Claude Clairaut（1713--1765），法国数学家，研究微分方程与天体力学。}它说明轨道只有在 \(r(s)\ge|j|\) 的区域内运动；非零角动量阻止单位速轨道到达旋转轴。

	固定 \(J=j\)，剩余的子午运动由一维有效势 \(V_j(s)=j^2/(2r(s)^2)\) 控制。设在 \([a,b]\) 的邻域内 \(r>0\)，函数
	\[
		f(s)=2E-\frac{j^2}{r(s)^2}
	\]
	在 \((a,b)\) 为正，在 \(a,b\) 有简单零点。曲线 \(p_s^2=f(s)\)、\(a\le s\le b\) 是一条光滑闭曲线：内部是上下两支图像，在端点用 \(p_s\) 作坐标，简单零点条件使 \(s\) 可由 \(p_s^2=f(s)\) 光滑解出。它与角圆的乘积就是一个紧正则共同水平集。取沿角圆的闭路和固定 \(\theta\) 沿子午闭曲线的闭路，作用分别为
	\begin{equation}\label{g4:eq:1104:revolution-actions}
		I_\theta=j,
		\qquad I_s(E,j)=\frac1\pi\int_a^b
		\sqrt{2E-\frac{j^2}{r(s)^2}}\,\mathrm{d}s.
	\end{equation}
	两支子午曲线对 \(\int p_s\,\mathrm{d}s\) 的贡献相等，所以出现因子 \(1/\pi\)。一个子午往返的时间与经度增量为
	\begin{equation}\label{g4:eq:1104:revolution-return}
		T_s=2\int_a^b\frac{\mathrm{d}s}{\sqrt{2E-j^2/r(s)^2}},
		\qquad
		\Delta\theta=2\int_a^b
		\frac{j\,\mathrm{d}s}{r(s)^2\sqrt{2E-j^2/r(s)^2}}.
	\end{equation}
	简单转点使两个积分的端点奇性至多为 \(|s-a|^{-1/2}\) 或 \(|b-s|^{-1/2}\)，因而有限。实际轨道闭合当且仅当 \(\Delta\theta/(2\pi)\) 为有理数：返回原状态必须经过整数次子午往返，而此时角增量是相应整数倍的 \(\Delta\theta\)。
\end{example}

平坦环面的两个平移与旋转曲面的一个旋转，都给出了可以直接计算的交换积分。对于一个 \(n\) 维黎曼流形，其测地流位于 \(2n\) 维余切丛；应用本章定理需要 \(n\) 个在稠密集上独立的对合积分。上述旋转曲面的计算完成了 \(n=2\) 时的独立性核验，更多维的对称性须逐一检查积分数目与交换关系。

\subsection{球摆及奇异纤维的入口}
\label{g4:subsec:1104:03}
球摆的摆杆长、质量和重力加速度均取为 \(1\)，竖直向上为第三坐标正向。位置 \(q\) 在单位球面上，利用球面的诱导度量把余切向量表示成满足 \(q\cdot p=0\) 的环境向量 \(p\)。于是相空间为
\[
	M=\{(q,p)\in\mathbb R^3\times\mathbb R^3:|q|=1,\ q\cdot p=0\}
	\cong\mathrm{T}^*S^2,
	\qquad \lambda=p\cdot\mathrm{d}q,\quad \omega=-\mathrm{d}\lambda.
\]
这个识别保持典范一形式：\(p\) 对底空间切向量的求值就是 Euclid 内积。因此 \(\omega\) 也等于环境二形式 \(\sum_i\mathrm{d}q_i\wedge\mathrm{d}p_i\) 在 \(M\) 上的限制。能量与竖直角动量为
\begin{equation}\label{g4:eq:1104:spherical-pendulum}
	H(q,p)=\frac{|p|^2}{2}+q_3,
	\qquad J(q,p)=(q\times p)_3,
	\qquad F=(H,J):M\longrightarrow\mathbb R^2.
\end{equation}
球摆及其积分的这一描述见 Cannas da Silva 第5.5节；下面把交换性、临界点和紧纤维逐项算出。\cite[第151--152页]{Cannas2006Symplectic}

\begin{proposition}\label{g4:prop:1104:pendulum-critical}
	式\eqref{g4:eq:1104:spherical-pendulum}中的 \(H,J\) Poisson 对合，其微分在稠密开集上线性无关。\(F\) 的秩零点恰为 \((q,p)=(e_3,0),(-e_3,0)\)，其值分别为 \((1,0),(-1,0)\)。其余临界点的底位置均在南半球的非极点处；写 \(q_3=c\in(-1,0)\)，每个这样的底位置恰有两个临界余切向量，且它们的值满足
	\begin{equation}\label{g4:eq:1104:pendulum-critical-values}
		j=\pm\frac{1-c^2}{\sqrt{-c}},
		\qquad h=\frac{3c}{2}-\frac1{2c}.
	\end{equation}
	每个非空正则共同水平集都是紧的，其每个连通分支具有\cref{g4:thm:1102:action-angle}给出的作用角邻域。
\end{proposition}
\begin{proof}
	先在不依赖球坐标的表示中计算向量场。点 \((q,p)\) 处的切向量 \((a,b)\) 满足 \(q\cdot a=0\)、\(q\cdot b+p\cdot a=0\)。直接代入约束导数可见下列两个向量场均沿 \(M\) 切向：
	\begin{equation}\label{g4:eq:1104:pendulum-fields}
		X_H=\bigl(p,-e_3+(q_3-|p|^2)q\bigr),
		\qquad X_J=(e_3\times q,e_3\times p).
	\end{equation}
	它们与 \((a,b)\) 的辛配对分别为
	\[
		p\cdot b+e_3\cdot a=\mathrm{d}H(a,b),\qquad
		(e_3\times q)\cdot b-(e_3\times p)\cdot a
		=e_3\cdot(q\times b+a\times p)=\mathrm{d}J(a,b).
	\]
	这证明了向量场公式。特别地，\(X_J\) 同时旋转 \(q,p\)，并且
	\[
		\{H,J\}=X_JH=p\cdot(e_3\times p)+e_3\cdot(e_3\times q)=0.
	\]

	离开两极，取余纬 \(\varphi\in(0,\pi)\) 与经度 \(\theta\)，使
	\[
		q=(\sin\varphi\cos\theta,\sin\varphi\sin\theta,\cos\varphi),
		\qquad \lambda=\eta\,\mathrm{d}\varphi+j\,\mathrm{d}\theta.
	\]
	度量为 \(\mathrm{d}\varphi^2+\sin^2\varphi\,\mathrm{d}\theta^2\)，从而
	\begin{equation}\label{g4:eq:1104:pendulum-spherical-coordinates}
		H=\frac12\left(\eta^2+\frac{j^2}{\sin^2\varphi}\right)+\cos\varphi,
		\qquad J=j.
	\end{equation}
	由于 \(\mathrm{d}J=\mathrm{d}j\ne0\)，秩下降恰好要求 \(H\) 对 \(\eta,\varphi\) 的两个导数都为零，即
	\[
		\eta=0,
		\qquad -\frac{j^2\cos\varphi}{\sin^3\varphi}-\sin\varphi=0.
	\]
	后一等式等价于 \(j^2\cos\varphi=-\sin^4\varphi\)，故 \(c=\cos\varphi<0\)，并给出式\eqref{g4:eq:1104:pendulum-critical-values}。每个底位置的两个符号对应反向的两个纬线运动，其角速度为
	\[
		\dot\theta=\frac{j}{\sin^2\varphi}=\pm\frac1{\sqrt{-c}}.
	\]
	这些点的秩为一。在非极点区域，\(\eta\ne0\) 已保证秩为二，因此正则点稠密。

	在两极，若 \(p\ne0\)，则 \(X_H\) 的底分量为非零向量 \(p\)，而 \(X_J=(0,e_3\times p)\ne0\)，两向量线性无关，故 \(\mathrm{d}H,\mathrm{d}J\) 仍独立。若 \(p=0\)，两个向量场都为零，得到所列两个秩零点。每个极点邻域也包含非极点的稠密正则点，所以整个相空间的正则集稠密。

	最后，\(H\ge |p|^2/2-1\)。对每个有限 \(E\)，次水平集 \(H\le E\) 是 \(\mathbb R^6\) 中的闭有界集，因而紧。于是任意纤维 \(F^{-1}(h,j)\) 都是紧集。正则纤维是无边界二维子流形，每个连通分支仍紧；两 Hamilton 场因 Poisson 对合而沿纤维切向。\cref{g4:thm:1102:action-angle}的独立性和紧连通性条件全部满足，故每个分支是 Lagrange 环面，并有该定理所述的邻域坐标。
\end{proof}

南半球的临界纬线是重力与绕轴转动共同形成的相对平衡。两个极点则是相空间中的真实平衡，它们的二次近似具有不同性质。

\begin{proposition}\label{g4:prop:1104:pendulum-linear-types}
	球摆北极平衡处，二次矩 \((H-1,J)\) 经实辛线性变换成为
	\begin{equation}\label{g4:eq:1104:pendulum-focus-pair}
		Q_1=u_1v_1+u_2v_2,
		\qquad Q_2=u_1v_2-u_2v_1,
		\qquad \omega_0=\mathrm{d}u_1\wedge\mathrm{d}v_1+\mathrm{d}u_2\wedge\mathrm{d}v_2.
	\end{equation}
	这是\secref{g4:sec:1103}中的焦点型二次矩。南极平衡处，二次矩经辛线性变换成为 \((I_++I_-,I_+-I_-)\)，其中 \(I_\pm=(Q_\pm^2+P_\pm^2)/2\)，即两个椭圆振子分量。
\end{proposition}
\begin{proof}
	在北极或南极附近用 \(x=q_1,y=q_2\) 作底坐标，写 \(q_3=\varepsilon\sqrt{1-x^2-y^2}\)，其中 \(\varepsilon=1\) 或 \(-1\)。以 \((\xi,\eta)\) 表示对应的典范余切坐标，记 \(x_0=(x,y)\)、\(\pi=(\xi,\eta)\)。这里 \(\pi\) 是 \(\lambda=\xi\,\mathrm{d}x+\eta\,\mathrm{d}y\) 的系数，须与环境三维向量 \(p\) 区分。局部度量矩阵及其逆为
	\[
		g=\operatorname{Id}+\frac{x_0x_0^{\mathsf T}}{1-|x_0|^2},
		\qquad g^{-1}=\operatorname{Id}-x_0x_0^{\mathsf T}.
	\]
	直接相乘可以核验逆矩阵。旋转的底向量场为 \(-y\partial_x+x\partial_y\)，故角动量精确等于 \(J=x\eta-y\xi\)。能量则为
	\[
		H=\frac12\bigl(|\pi|^2-(x_0\cdot\pi)^2\bigr)
		+\varepsilon\sqrt{1-|x_0|^2}
		=\varepsilon+\frac12\bigl(|\pi|^2-\varepsilon|x_0|^2\bigr)+O(|(x_0,\pi)|^4).
	\]
	在北极取
	\[
		u=\frac{x_0+\pi}{\sqrt2},\qquad
		v=\frac{\pi-x_0}{\sqrt2}.
	\]
	逐分量计算得到 \(\mathrm{d}u_j\wedge\mathrm{d}v_j=\mathrm{d}(x_0)_j\wedge\mathrm{d}\pi_j\)，并有
	\[
		\frac{|\pi|^2-|x_0|^2}{2}=u\cdot v,
		\qquad x\eta-y\xi=u_1v_2-u_2v_1.
	\]
	因此得到所述焦点型二次矩。若 \(R=\left(\begin{smallmatrix}0&-1\\1&0\end{smallmatrix}\right)\)，则
	\[
		X_{aQ_1+bQ_2}(u,v)=(au+bRu,-av+bRv).
	\]
	当 \(ab\ne0\) 时，其复特征值为 \(a\pm\mathrm{i}b,-a\pm\mathrm{i}b\)。特别地，原能量的二次部分 \(Q_1\) 自身给出实指数伸缩；焦点型描述的是两个交换二次 Hamilton 函数组成的矩。

	在南极，能量二次部分为 \(H_2=(x^2+y^2+\xi^2+\eta^2)/2\)。令
	\[
		Q_+=\frac{x+\eta}{\sqrt2},\quad
		P_+=\frac{\xi-y}{\sqrt2},\qquad
		Q_-=\frac{x-\eta}{\sqrt2},\quad
		P_-=\frac{\xi+y}{\sqrt2}.
	\]
	相加外积得到 \(\mathrm{d}Q_+\wedge\mathrm{d}P_++\mathrm{d}Q_-\wedge\mathrm{d}P_-=\mathrm{d}x\wedge\mathrm{d}\xi+\mathrm{d}y\wedge\mathrm{d}\eta\)。展开平方又得 \(H_2=I_++I_-\)、\(J=I_+-I_-\)，故是两个独立椭圆分量。此外，完整能量满足 \(H+1=|p|^2/2+(q_3+1)\ge0\)，等号只在南极零动量成立，这也直接显示了该点的严格能量极小性质。
\end{proof}

球摆的上方平衡所在纤维还可以整体参数化，无需从二次近似推断它的拓扑。

\begin{example}\label{g4:exm:1104:pendulum-pinched-torus}
	考虑奇异纤维 \(F^{-1}(1,0)\)。在闭圆柱 \([0,2\pi]\times(\mathbb R/2\pi\mathbb Z)\) 上令
	\begin{equation}\label{g4:eq:1104:pendulum-singular-parametrization}
		\begin{aligned}
			q(\tau,\alpha)&=(\sin\tau\cos\alpha,\sin\tau\sin\alpha,\cos\tau),\\
			p(\tau,\alpha)&=2\sin(\tau/2)\,\partial_\tau q(\tau,\alpha).
		\end{aligned}
	\end{equation}
	因为 \(|\partial_\tau q|=1\) 且它与 \(q\) 正交，这些点属于 \(M\)，并满足
	\[
		H=2\sin^2(\tau/2)+\cos\tau=1,\qquad J=0.
	\]
	圆柱两条边界均映到同一个点 \((e_3,0)\)。内部则与其余纤维双射。事实上，非极点处 \(J=0\) 使 \(p\) 沿子午方向，能量决定 \(|p|=\sqrt{2(1-q_3)}\)；在给定 \(q\) 时，参数 \(\tau\) 与 \(2\pi-\tau\)、经度相差 \(\pi\) 的两种写法给出相反的 \(p\)，所以动量方向唯一确定所用写法。在南极 \(\tau=\pi\)，\(p=-2(\cos\alpha,\sin\alpha,0)\)，也恰好一次遍历允许的动量圆。在北极，能量方程强制 \(p=0\)。

	把圆柱的两条边界圆全部压成同一个点，所得紧商空间因而连续双射到 \(F^{-1}(1,0)\)。后者是 Hausdorff 空间中的闭子集，所以这是一张同胚。这个商也可以先把圆柱两边对应粘合成环面，再把粘合圆压成一点；因此球摆的这条实际纤维是针缩环面。其余点都是秩二点，奇异性集中在北极平衡处。沿式\eqref{g4:eq:1104:pendulum-singular-parametrization}中固定 \(\alpha\) 的运动，\(\dot\tau=2\sin(\tau/2)\)；到达任一端点所需时间是发散的积分 \(\int\mathrm{d}\tau/(2\sin(\tau/2))\)，所以这些非平衡轨道在两个时间方向上都渐近北极。
\end{example}

由此，球摆在正则值处给出有作用角邻域的紧环面，在 \((1,0)\) 处给出实际的针缩纤维，并在其奇点具有与前节代表模型一致的焦点型二次矩。作用角坐标的局部构造、临界纤维的拓扑及沿基空间闭路的周期格变化，是三个可以分别计算的层次；前节的代表模型已展示最后一种变化怎样产生整体单值障碍。
% 纲要条目排为编号小节标题，并进入目录与书签；说明段仍保留于纲要注释中。


\begin{chaptersummary}
在正则共同水平集上，独立对合积分的 Hamilton 场既切于水平集，又张成其全部切空间，因而纤维是 Lagrange 的。完备性交换流给出传递的 \(\mathbb R^n\) 作用，离散稳定子记录全部返回时间；紧连通性使该稳定子成为满秩格，纤维于是为环面。

\ProseParagraphBreak

紧纤维附近可以延拓一组完整的周期基。沿其基本闭路积分原始一形式并除以 \(2\pi\)，所得作用变量的微分恰好归一化周期向量场。局部截面确定角的原点，再以底空间上的一形式调整角变量，就得到 \(\omega=\sum_j\mathrm{d}\theta_j\wedge\mathrm{d}I_j\)。在这些坐标中，\(H=h(I)\) 的运动为固定作用环面上的匀速平移，频率向量是 \(\nabla h\)。

\ProseParagraphBreak

局部作用变量的换基属于整数仿射变换；其微分组成周期格局部系统。沿底空间环路延拓后，基本闭路可以发生非平凡整数换基，从而阻止全局作用角坐标。显式焦点型模型的第二个周期包含 \(\arg c\)，绕奇值一周后增加第一个周期；它的零纤维则是针缩环面。即使线性单值平凡，作用的一形式周期和角原点的相容性仍须另外检查。

\ProseParagraphBreak

振子直接展示作用、频率与共振，旋转曲面的测地流把角动量守恒化为一维有效势问题，球摆同时包含正则紧环面、相对平衡和针缩奇纤维。可积性规定的是足够多的独立交换积分；它与某条轨道是否闭合、某个积分选择是否在奇点退化，是需要分别计算的性质。
\end{chaptersummary}

\BookExercises
\begin{exercise}\label{g4:ex:11:compact-component-only}
	在标准辛平面 \((\mathbb R^2_{q,p},\mathrm{d}q\wedge\mathrm{d}p)\) 上取
	\[
		H(q,p)=(q^2+p^2-1)(q-2).
	\]
	证明 \(0\) 为 \(H\) 的正则值，并确定 \(H^{-1}(0)\) 的全部连通分支。在紧分支上计算 Hamilton 流的最小正周期；在非紧分支上求出过任意初值的最大积分曲线及其时间区间。说明该例如何区分所选纤维紧、整个水平集紧以及环境向量场完备这三个条件。
\end{exercise}
\begin{solution}
	零水平集为两个不相交的连通集合
	\[
		L_{\mathrm c}=\{q^2+p^2=1\},\qquad
		L_{\mathrm l}=\{q=2\}.
	\]
	它们分别为圆周与直线，且彼此有正距离，所以正是两个连通分支。微分为
	\[
		\mathrm{d}H=
		\bigl(2q(q-2)+q^2+p^2-1\bigr)\,\mathrm{d}q
		+2p(q-2)\,\mathrm{d}p.
	\]
	在圆周上，它等于 \(2(q-2)(q\,\mathrm{d}q+p\,\mathrm{d}p)\)，处处非零；在直线上，它等于 \((3+p^2)\,\mathrm{d}q\)，也处处非零。因此 \(0\) 是正则值。

	Hamilton 方程为
	\[
		\dot q=2p(q-2),\qquad
		\dot p=-2q(q-2)-(q^2+p^2-1).
	\]
	在 \(L_{\mathrm c}\) 上写 \(q=\cos\varphi\)、\(p=\sin\varphi\)，得到
	\[
		\dot\varphi=2(2-\cos\varphi)>0.
	\]
	角度严格递增，所以第一次回到起点恰是绕行一周，最小正周期与起点无关，等于
	\[
		T=\int_{-\pi}^{\pi}\frac{\mathrm{d}\varphi}{2(2-\cos\varphi)}
		=\int_{-\infty}^{\infty}\frac{\mathrm{d}u}{1+3u^2}
		=\frac{\pi}{\sqrt3}.
	\]
	中间用了 \(u=\tan(\varphi/2)\)、\(\cos\varphi=(1-u^2)/(1+u^2)\)。紧纤维上的场完备，与这个周期运动的显式计算一致。

	在 \(L_{\mathrm l}\) 上，\(q(t)=2\)，而 \(\dot p=-(3+p^2)\)。令 \(a=\arctan(p_0/\sqrt3)\in(-\pi/2,\pi/2)\)，解为
	\[
		p(t)=\sqrt3\tan(a-\sqrt3\,t),\qquad
		\frac{a-\pi/2}{\sqrt3}<t<\frac{a+\pi/2}{\sqrt3}.
	\]
	在两个有限端点处 \(|p(t)|\to\infty\)，所以这个区间就是最大存在区间，向量场在该纤维上及在环境 \(\mathbb R^2\) 上都不完备。整个零水平集也不紧，但其中的 \(L_{\mathrm c}\) 满足 \cref{g4:thm:1101:liouville-torus} 的全部条件；定理只对所选组件要求紧性。
\end{solution}

\begin{exercise}\label{g4:ex:11:singular-cross-fibre}
	在标准辛平面上取 \(H(q,p)=(qp)^2\)。证明它给出一个一自由度可积系统，求全部正水平集的连通分支及其 Hamilton 流，并判定这些受限流是否完备。再说明零纤维不在原点附近构成一维流形，并求零纤维上每一点的运动。
\end{exercise}
\begin{solution}
	有
	\[
		\mathrm{d}H=2qp^2\,\mathrm{d}q+2q^2p\,\mathrm{d}p,
		\qquad
		X_H=2q^2p\,\partial_q-2qp^2\,\partial_p.
	\]
	微分非零当且仅当 \(qp\neq0\)，这是稠密开集；一个函数与自身的 Poisson 括号总为零，所以满足一自由度可积性的定义。对 \(c>0\)，方程 \((qp)^2=c\) 等价于 \(qp=k\)，其中 \(k=\sqrt c\) 或 \(-\sqrt c\)。对每个 \(k\)，按 \(q>0\) 或 \(q<0\) 又分成两支，合计四个正则连通分支。

	在其中一支上，\(p=k/q\)，Hamilton 方程化为
	\[
		\dot q=2kq,\qquad \dot p=-2kp,
	\]
	所以过 \((q_0,p_0)\) 的解为
	\[
		q(t)=q_0e^{2kt},\qquad p(t)=p_0e^{-2kt}.
	\]
	解对全部 \(t\in\mathbb R\) 存在，始终留在原来的分支上。用 \(u=\log|q|\) 参数化该分支，得到 \(u(t)=u(0)+2kt\)，因此它微分同胚于 \(\mathbb R\)，受限流完备且没有非零周期。这对应 \cref{g4:prop:1101:complete-fibre-cylinder} 中 \(n=1,r=0\) 的情形。

	零纤维为 \(\{q=0\}\cup\{p=0\}\)，即两条坐标轴的并。在原点的任意充分小圆盘中，移去原点后有四条连通的半轴；若存在一维流形图，可再缩小到映为开区间的邻域，移去中心点后至多只有两个连通分支。由于一个包含原点的相对开邻域必定包含四条半轴的初始小段，而它们在整个零纤维中只能通过原点相连，矛盾。故零纤维在原点不是一维流形。这里 \(\mathrm{d}H\) 在整个零纤维上都为零，相应 \(X_H\) 也为零，所以每一点都是静止轨道。纤维失去正则性后，交换场不能再构成其切向标架。
\end{solution}

\begin{exercise}\label{g4:ex:1102:nonlinear-integral-parameters}
	在 \(\mathrm{T}^*\mathbb T^2\) 上取
	\[
		\lambda=p_1\,\mathrm{d}q_1+p_2\,\mathrm{d}q_2,\qquad
		\omega=-\mathrm{d}\lambda,\qquad
		q_i\in\mathbb R/(2\pi\mathbb Z),
	\]
	并令 \(F=(f_1,f_2)=(p_1,e^{p_1}p_2)\)。证明 \(F\) 的每个共同水平集都是紧正则环面。以 \(q_1,q_2\) 的正向基本圈为基，求原场 \((X_{f_1},X_{f_2})\) 的完整周期格及作用变量，核对作用微分公式。再对 \(H=f_2\) 求环面上的频率，并判定轨道何时闭合、闭合时的最小正周期是多少。
\end{exercise}
\begin{solution}
	写 \(F=b=(b_1,b_2)\)，则
	\[
		p_1=b_1,\qquad p_2=e^{-b_1}b_2.
	\]
	映射 \(p\mapsto(p_1,e^{p_1}p_2)\) 的 Jacobian 行列式为 \(e^{p_1}>0\)，且上述公式给出整体光滑逆。两分量只依赖动量，因而 Poisson 对合。每张纤维是这组固定动量上的完整 \(\mathbb T^2\)，紧且正则。

	由 Hamilton 方程，
	\[
		X_{f_1}=\partial_{q_1},\qquad
		X_{f_2}=e^{p_1}p_2\,\partial_{q_1}
			+e^{p_1}\partial_{q_2}.
	\]
	在 \(F=b\) 的纤维上，多时间 \(t=(t_1,t_2)\) 使角变量变化
	\[
		\Delta q_1=t_1+b_2t_2,\qquad
		\Delta q_2=e^{b_1}t_2.
	\]
	回到原点当且仅当这两个增量属于 \(2\pi\mathbb Z\)。解出 \(t\)，得到完整周期格
	\[
		\Gamma_b=\Lambda(b)\mathbb Z^2,\qquad
		\Lambda(b)=2\pi
		\begin{pmatrix}
			1&-b_2e^{-b_1}\\
			0&e^{-b_1}
		\end{pmatrix}.
	\]
	第一列对应 \(q_1\) 正向一周；第二列使 \(q_1\) 不变而 \(q_2\) 正向一周。上述“当且仅当”也排除了额外周期。

	沿这两个基本圈积分 \(\lambda\)，得到
	\[
		I_1=b_1,\qquad I_2=e^{-b_1}b_2.
	\]
	因此
	\[
		\begin{pmatrix}\mathrm{d}I_1\\ \mathrm{d}I_2\end{pmatrix}
		=
		\begin{pmatrix}
			1&0\\
			-b_2e^{-b_1}&e^{-b_1}
		\end{pmatrix}
		\begin{pmatrix}\mathrm{d}b_1\\ \mathrm{d}b_2\end{pmatrix}
		=\frac1{2\pi}\Lambda(b)^{\mathsf T}\mathrm{d}b.
	\]
	可直接取作用角坐标 \((q_1,q_2,I_1,I_2)\)，其作用就是原来的动量。虽然 \(F\) 是另一组全局纤维参数，它不是这组基本圈对应的作用变量。

	对 \(H=e^{I_1}I_2\)，频率为
	\[
		(\dot q_1,\dot q_2)=(e^{I_1}I_2,e^{I_1}).
	\]
	第二个频率始终非零。若正时间 \(T\) 使轨道闭合，必有
	\(e^{I_1}T=2\pi k\)、\(I_2e^{I_1}T=2\pi m\)，其中 \(k\) 为正整数、\(m\in\mathbb Z\)，故 \(I_2=m/k\) 有理。反之，若 \(I_2=m/k\) 为既约分数，\(k>0\)，则
	\[
		T_{\min}=2\pi k\,e^{-I_1}.
	\]
	既约性保证没有更小的正整数 \(k\) 满足两个闭合条件。特别地，\(I_2=0\) 对应 \(m=0,k=1\)，仍为非静止闭轨道。
\end{solution}

\begin{exercise}\label{g4:ex:1102:power-well-action}
	在 \((\mathbb R^2_{q,p},\mathrm{d}q\wedge\mathrm{d}p)\) 上取
	\[
		H(q,p)=\frac{p^2}{2}+\frac{q^{2m}}{2m},
		\qquad m\ge2\text{ 为整数}.
	\]
	对 \(E>0\)，沿 Hamilton 运动方向为 \(H^{-1}(E)\) 定向，并取原始形式 \(\lambda=p\,\mathrm{d}q\)。证明该水平集是一条紧正则闭曲线，计算其作用 \(I(E)\)、周期 \(T(E)\) 及作用角频率。把常数写成 \([0,1]\) 上的定积分即可。另由 \(\dot q=p\) 直接求周期，核对 \(T(E)=2\pi I'(E)\)，并说明 \(E\downarrow0\) 时的行为。
\end{exercise}
\begin{solution}
	微分
	\(\mathrm{d}H=q^{2m-1}\mathrm{d}q+p\,\mathrm{d}p\)
	只在原点为零，故每个 \(E>0\) 正则。水平集由上下两条图
	\[
		p=\pm\sqrt{2E-\frac{q^{2m}}m},
		\qquad -a\le q\le a,\qquad a=(2mE)^{1/(2m)}
	\]
	在两个端点相接而成，闭且有界。端点附近可把 \(q\) 光滑地解为 \(p\) 的函数，因为那里 \(q^{2m-1}\ne0\)；所以这确是一条光滑紧连通闭曲线。

	上半支 \(p>0\) 时 \(\dot q=p>0\)，运动从 \(-a\) 到 \(a\)；下半支则反向返回。两支的 \(p\,\mathrm{d}q\) 积分均为正。因此
	\[
		\begin{aligned}
			I(E)
			&=\frac1\pi\int_{-a}^a
			\sqrt{2E-\frac{q^{2m}}m}\,\mathrm{d}q\\
			&=C_m E^{(m+1)/(2m)},\\
			C_m&=\frac{2\sqrt2}{\pi}(2m)^{1/(2m)}
			\int_0^1\sqrt{1-u^{2m}}\,\mathrm{d}u.
		\end{aligned}
	\]
	于是 \(H\) 在作用变量中为
	\[
		H(I)=\left(\frac I{C_m}\right)^{2m/(m+1)},\qquad
		\dot\theta=\frac{2m}{m+1}\,
			C_m^{-2m/(m+1)}I^{(m-1)/(m+1)}.
	\]
	角的周期为 \(2\pi\)，所以
	\[
		T(E)=2\pi I'(E)
		=2\pi C_m\frac{m+1}{2m}\,
			E^{(1-m)/(2m)}.
	\]

	现在直接按时间积分核验。上下两支各走一次，给出
	\[
		T(E)=2\int_{-a}^a
		\frac{\mathrm{d}q}{\sqrt{2E-q^{2m}/m}}
		=\frac{4a}{\sqrt{2E}}
		\int_0^1\frac{\mathrm{d}u}{\sqrt{1-u^{2m}}}.
	\]
	端点的被积函数只有可积的平方根奇性，积分有限。令
	\[
		J_m=\int_0^1\sqrt{1-u^{2m}}\,\mathrm{d}u,\qquad
		K_m=\int_0^1(1-u^{2m})^{-1/2}\,\mathrm{d}u.
	\]
	分部积分得到
	\[
		J_m
		=m\int_0^1\frac{u^{2m}}{\sqrt{1-u^{2m}}}\,\mathrm{d}u
		=m(K_m-J_m),
		\qquad J_m=\frac m{m+1}K_m.
	\]
	可先在 \([0,1-\varepsilon]\) 分部积分再令 \(\varepsilon\downarrow0\)，边界项 \(u\sqrt{1-u^{2m}}\) 趋于零。将这个关系代入 \(C_m\)，两种周期表达完全一致。

	当 \(E\downarrow0\) 时，\(I(E)\to0\)、\(T(E)\to+\infty\)，而角频率趋于零。零水平缩成唯一临界点，已不是正则闭曲线；作用角邻域只围绕各个 \(E>0\) 的环面建立，不能把极限点也视为其中一张正则纤维。
\end{solution}

\begin{exercise}\label{g4:ex:1103:monodromy-cover}
	在\cref{g4:prop:1103:focus-monodromy}的模型中，计算绕奇值正向 \(k\) 周的单值，其中 \(k\in\mathbb Z\)。求保持不变的整数同调类。若将穿孔底圆盘拉回到它的一个连通有限叶覆盖上，单值障碍是否消失？若拉回到万有覆盖，能否因此断言已经获得全局作用角坐标？
\end{exercise}
\begin{solution}
	写 \(C=I+N\)，其中 \(N=\left(\begin{smallmatrix}0&1\\0&0\end{smallmatrix}\right)\)、\(N^2=0\)。直接相乘及求逆得
	\[
		C^k=I+kN=\begin{pmatrix}1&k\\0&1\end{pmatrix}
	\]
	对所有整数 \(k\) 都成立。类 \(a\gamma_1+b\gamma_2\) 经过一次正向延拓后变成 \((a+b)\gamma_1+b\gamma_2\)，故保持不变的类恰为 \(\mathbb Z\gamma_1\)。

	对任意连通 \(d\) 叶覆盖，固定基点上方一点。一次正向底环路在这个有限纤维上诱导置换，因而存在某个正整数 \(k\) 使走 \(k\) 周后返回所选点。此时提升为覆盖空间中的闭路，其拉回单值为 \(C^k\ne I\)。所以任何有限叶覆盖都保留非平凡单值。

	在万有覆盖上，每条闭路都零伦，故周期格单值消失。对本模型，\(\arg c\) 有全局实提升，式\eqref{g4:eq:1103:focus-actions}也成为全局实值函数。但从“单值消失”这个事实本身只能推出周期格的全局基；一般还要核对闭一形式的恰当性、作用映射的全局单射性及角原点的相容性，不能把它直接当作全局作用角定理。\cref{g4:exm:1103:affine-circle,g4:exm:1103:magnetic-obstruction}分别说明其中的额外问题。
\end{solution}

\begin{exercise}\label{g4:ex:1103:focus-frequency}
	在焦点粘合模型的一个辐角分支上记 \(c=a+\mathrm{i}b\)、\(\tau=-\log|c|\)、\(\varphi=\arg c\)，采用式\eqref{g4:eq:1103:focus-actions}的作用变量。求 Hamilton 函数 \(H=a\) 在相应角坐标中的频率。证明这个流在给定正则环面上闭合当且仅当 \(\varphi/(2\pi)\) 为有理数；若该数为既约分数 \(p/q\)，其中 \(q>0\)，求最小正周期，并检验改变辐角分支不会改变结论。
\end{exercise}
\begin{solution}
	由 \(\mathrm{d}I_1=\mathrm{d}b\) 和 \(2\pi\,\mathrm{d}I_2=\tau\,\mathrm{d}a+\varphi\,\mathrm{d}b\)，得到
	\[
		\mathrm{d}H=-\frac{\varphi}{\tau}\,\mathrm{d}I_1
		+\frac{2\pi}{\tau}\,\mathrm{d}I_2,
		\qquad \nu=\left(-\frac{\varphi}{\tau},\frac{2\pi}{\tau}\right).
	\]
	时间 \(T>0\) 返回要求 \(T/\tau\) 是正整数 \(m\)，并且 \(m\varphi\in2\pi\mathbb Z\)。故闭合恰等价于所述有理条件；若 \(\varphi/(2\pi)=p/q\) 既约，则最小 \(m=q\)，周期为 \(q\tau\)。此结论也可直接由 \((T,0)\) 是否落入时间周期格检验。

	改变分支使 \(\varphi\) 增加 \(2\pi k\)，有理性及既约分母 \(q\) 不变，周期也不变。沿固定有理辐角令 \(|c|\downarrow0\)，有 \(T=-q\log|c|\to\infty\)；这显示了接近奇纤维时的对数时间增长。
\end{solution}

\begin{exercise}\label{g4:ex:1103:magnetic-section}
	设 \(n\ge2\)，\(B\subset\mathbb R^n\) 为非空连通开集，底坐标为 \(b_1,\ldots,b_n\)，\(\beta\) 为 \(B\) 上的光滑闭二形式。记 \(\mathbb T^n=(\mathbb R/2\pi\mathbb Z)^n\)，在 \(M=\mathbb T^n\times B\) 上设
	\[
		\omega_0=\sum_{j=1}^n\mathrm{d}\theta_j\wedge\mathrm{d}b_j,
		\qquad \omega_\beta=\omega_0+\pi^*\beta,
		\qquad \pi(\theta,b)=b.
	\]
	\begin{enumerate}
		\item 设 \(\beta=\mathrm{d}\alpha\)，其中 \(\alpha=\sum_j a_j(b)\,\mathrm{d}b_j\) 是 \(B\) 上的光滑一形式。证明
		\[
			\Phi_\alpha(\theta,b)=(\theta+a(b),b)
		\]
		是从 \((M,\omega_\beta)\) 到 \((M,\omega_0)\) 的全局辛同胚。求出它所给的作用角坐标及一个 Lagrange 截面。
		\item 若 \(\alpha'\) 是 \(B\) 上 \(\beta\) 的另一原始一形式，记 \(\eta=\alpha'-\alpha=\sum_j c_j(b)\,\mathrm{d}b_j\)。计算 \(\Phi_{\alpha'}\Phi_\alpha^{-1}\)，并证明这个变换是标准辛结构下某个仅依赖 \(b\) 的自主 Hamilton 函数的时间 \(1\) 流，当且仅当 \(\eta\) 恰当。
		\item 在 \(n=2\)、\(B=\mathbb R^2\setminus\{0\}\) 时，取 \(\beta=0\)、\(\alpha=0\) 及
		\[
			\alpha'=\kappa\frac{-b_2\,\mathrm{d}b_1+b_1\,\mathrm{d}b_2}{b_1^2+b_2^2},
			\qquad \kappa\in\mathbb R\setminus\{0\}.
		\]
		写出对应的全局辛角平移，并证明它不能由第二问所述的 Hamilton 函数实现。
	\end{enumerate}
\end{exercise}
\begin{solution}
	1. 每个 \(a_j\) 都是全局实值光滑函数，因此加上 \(a(b)\) 在角圆上定义良好，逆映射为 \((\theta,b)\mapsto(\theta-a(b),b)\)。直接计算
	\[
		\Phi_\alpha^*\omega_0
		=\sum_j\bigl(\mathrm{d}\theta_j+\mathrm{d}a_j\bigr)\wedge\mathrm{d}b_j
		=\omega_0+\pi^*(\mathrm{d}\alpha)=\omega_\beta.
	\]
	故它是所述辛同胚，作用角坐标可取 \(I_j=b_j\)、\(\vartheta_j=\theta_j+a_j(b)\)。把修正角固定为零，得到截面 \(s(b)=(-a(b),b)\)，且
	\(s^*\omega_\beta=-\mathrm{d}\alpha+\beta=0\)，所以这个截面是 Lagrange 的。

	2. 因 \(\mathrm{d}\eta=\mathrm{d}\alpha'-\mathrm{d}\alpha=0\)，两套坐标间的变换为
	\[
		T_c(\theta,b)=(\theta+c(b),b),
		\qquad T_c^*\omega_0=\omega_0+\pi^*(\mathrm{d}\eta)=\omega_0.
	\]
	对仅依赖 \(b\) 的函数 \(H\)，本章的符号约定给出
	\[
		X_H=\sum_j\frac{\partial H}{\partial b_j}\partial_{\theta_j},
		\qquad \varphi_H^t(\theta,b)=(\theta+t\nabla H(b),b).
	\]
	这个流对全部实时间存在。若 \(\eta=\mathrm{d}H\)，便有 \(T_c=\varphi_H^1\)。反过来，若 \(T_c=\varphi_H^1\)，由于角的周期是 \(2\pi\)，有
	\[
		c(b)-\nabla H(b)\in2\pi\mathbb Z^n.
	\]
	左端连续，而 \(B\) 连通，所以它等于某个常向量 \(2\pi k\)、\(k\in\mathbb Z^n\)。于是
	\[
		\eta=\mathrm{d}\!\left(H+2\pi\sum_jk_jb_j\right),
	\]
	仍然恰当。

	3. 在穿孔平面上直接求外微分得到 \(\mathrm{d}\alpha'=0\)，而对应的变换为
	\[
		c(b)=\frac{\kappa}{b_1^2+b_2^2}(-b_2,b_1),
		\qquad T_c(\theta,b)=(\theta+c(b),b),
	\]
	角分量均按 \(2\pi\) 取模。沿正向单位圆 \(b(t)=(\cos t,\sin t)\)、\(0\le t\le2\pi\)，有
	\[
		\int_{S^1}\alpha'=\int_0^{2\pi}\kappa\,\mathrm{d}t=2\pi\kappa\ne0.
	\]
	所以 \(\alpha'\) 不恰当。这个角平移保持标准辛形式，但依第二问，不能是任何仅依赖底变量的自主 Hamilton 函数的时间 \(1\) 流。
\end{solution}

\begin{exercise}\label{g4:ex:1104:resonant-polynomial}
	设 \(m,n\) 为互素正整数，\(\nu_0>0\)，在标准 \(\mathbb R^4\) 上取
	\[
		I_j=(q_j^2+p_j^2)/2,\qquad
		H=\nu_0(mI_1+nI_2),\qquad
		a_j=q_j-\mathrm{i}p_j.
	\]
	证明复多项式 \(K=a_1^n\overline{a_2}^{\,m}\) 的实部与虚部均为第一积分，并求 \(|K|^2\)。在 \(I_1,I_2>0\) 的固定环面上，证明每条轨道的最小正周期为 \(2\pi/\nu_0\)，其整数绕行数为 \((m,n)\)；再证明 \(K\) 的一个非空水平集在该环面上恰是一条轨道。
\end{exercise}
\begin{solution}
	Hamilton 方程给出 \(\dot q_1=m\nu_0p_1\)、\(\dot p_1=-m\nu_0q_1\)，故 \(\dot a_1=\mathrm{i}m\nu_0a_1\)；同理 \(\dot a_2=\mathrm{i}n\nu_0a_2\)。乘积法则于是给出
	\[
		\dot K
		=n\mathrm{i}m\nu_0K-m\mathrm{i}n\nu_0K=0,
		\qquad |K|^2=(2I_1)^n(2I_2)^m.
	\]
	所以实部与虚部都是全相空间上的实多项式第一积分。

	在正作用环面上，\(a_j=\sqrt{2I_j}\,e^{\mathrm{i}\theta_j}\)，角速度为 \((m\nu_0,n\nu_0)\)。若 \(T>0\) 是周期，则 \(m\nu_0T/(2\pi)\)、\(n\nu_0T/(2\pi)\) 均为整数。由 \(m,n\) 互素，存在整数 \(r,s\) 使 \(rm+sn=1\)，故 \(\nu_0T/(2\pi)\) 也是整数。最小正取值为 \(1\)，给出所述周期；相应角增量为 \(2\pi(m,n)\)。

	在固定环面上，\(K\) 的模长是非零常数，它的值只规定 \(n\theta_1-m\theta_2\) 模 \(2\pi\)。若两个点之差为 \((\delta_1,\delta_2)\)，且
	\[
		n\delta_1-m\delta_2=2\pi\ell,\qquad \ell\in\mathbb Z,
	\]
	则可取整数 \(a,b\) 使 \(na-mb=-\ell\)。于是 \(n(\delta_1+2\pi a)=m(\delta_2+2\pi b)\)，所以存在实数 \(t\) 使
	\[
		(\delta_1+2\pi a,\delta_2+2\pi b)=t\nu_0(m,n).
	\]
	这正表示两点在同一轨道上。反过来 \(K\) 沿轨道恒定，故非空水平集恰是一条轨道。这里互素性保证相位方程的水平集连通。
\end{solution}

\begin{exercise}\label{g4:ex:1104:latitude-stability}
	在度量 \(\mathrm{d}s^2+r(s)^2\mathrm{d}\theta^2\)、\(r>0\) 的旋转区域内，固定非零角动量 \(j\)。证明 \(s=s_0\) 上存在非静止纬线测地线当且仅当 \(r'(s_0)=0\)。求其周期，并在线性化的固定 \(j\) 子午方程中判定 \(r''(s_0)<0\)、\(r''(s_0)>0\) 两种情形。对单位速轨道，将小子午振动的角频率写成 \(r(s_0),r''(s_0)\) 的表达式。
\end{exercise}
\begin{solution}
	固定 \(j\) 后的方程为
	\[
		\dot s=p_s,\qquad
		\dot p_s=\frac{j^2r'(s)}{r(s)^3},\qquad
		\dot\theta=\frac{j}{r(s)^2}.
	\]
	要使 \(s=s_0\) 恒定，必须且只须 \(p_s=0\)、\(r'(s_0)=0\)。此时经度以非零常速度 \(j/r(s_0)^2\) 变化，最小正周期为
	\[
		T=\frac{2\pi r(s_0)^2}{|j|}.
	\]
	令 \(\sigma=\delta s\)、\(\rho=\delta p_s\)。在该相对平衡处线性化，得到
	\[
		\dot\sigma=\rho,\qquad
		\dot\rho=\frac{j^2r''(s_0)}{r(s_0)^3}\,\sigma.
	\]
	若 \(r''(s_0)<0\)，线性化子午运动为振子，角频率是
	\[
		\Omega_s=|j|\sqrt{\frac{-r''(s_0)}{r(s_0)^3}}.
	\]
	若 \(r''(s_0)>0\)，则出现一正一负的实特征值，存在指数增长的子午扰动。单位速纬线满足 \(|j|=r(s_0)\)，因而 \(T=2\pi r(s_0)\)、\(\Omega_s=\sqrt{-r''(s_0)/r(s_0)}\)。这些结论描述固定角动量下的横向线性化；若 \(r''(s_0)=0\)，二阶线性化不能据此确定稳定类型。例如单位球面以余纬弧长为 \(s\) 时 \(r(s)=\sin s\)，赤道处 \(r=1,r''=-1\)，故 \(\Omega_s=1\)。
\end{solution}

\begin{exercise}\label{g4:ex:1104:pendulum-period-limit}
	对球摆 \(H=|p|^2/2+q_3\)，证明每个通过竖直轴的平面与球面相交得到的不变子系统都可以写为
	\[
		q(\psi)=(\sin\psi\cos\alpha,\sin\psi\sin\alpha,-\cos\psi),
		\qquad H=\frac12\dot\psi^2-\cos\psi,
	\]
	其中 \(\alpha\) 固定且 \(J=0\)。对 \(-1<h<1\) 的振动轨道，证明其周期为
	\[
		T(h)=4\int_0^{\pi/2}\frac{\mathrm{d}\chi}{\sqrt{1-k^2\sin^2\chi}},
		\qquad k^2=\frac{h+1}{2}.
	\]
	由此证明 \(T(h)\to2\pi\) 当 \(h\downarrow-1\)，而 \(T(h)\to+\infty\) 当 \(h\uparrow1\)，并说明后一极限与上平衡纤维中的轨道相容。
\end{exercise}
\begin{solution}
	若初始 \(q,p\) 均在含竖直轴的固定平面 \(P\) 中，式\eqref{g4:eq:1104:pendulum-fields}的右侧仍在 \(P\times P\) 中，因为 \(e_3\in P\)。因此该约束子流形上的方程有解，唯一性使这些解同时是原方程的解，故它不变。用所给参数表示其位置，\(|\partial_\psi q|=1\)，并由 \(\dot q=p\) 得到 \(p=\dot\psi\,\partial_\psi q\)，从而能量如题中所示。\(q,p\) 共处竖直平面也给出 \((q\times p)_3=0\)。

	固定 \(-1<h<1\)。转点为 \(\psi=\pm\psi_*\)，其中 \(\psi_*\in(0,\pi)\)、\(\cos\psi_*=-h\)。一个完整周期由四个对称的四分之一振动组成，所以
	\[
		T(h)=4\int_0^{\psi_*}\frac{\mathrm{d}\psi}{\sqrt{2(h+\cos\psi)}}.
	\]
	令 \(\sin(\psi/2)=k\sin\chi\)，其中 \(k=\sin(\psi_*/2)=\sqrt{(h+1)/2}\)。当 \(\psi\) 从 \(0\) 增加到 \(\psi_*\) 时，\(\chi\) 从 \(0\) 增加到 \(\pi/2\)，且
	\[
		\mathrm{d}\psi=\frac{2k\cos\chi}{\sqrt{1-k^2\sin^2\chi}}\,\mathrm{d}\chi,
		\qquad \sqrt{2(h+\cos\psi)}=2k\cos\chi.
	\]
	消去公共因子即得所述积分。

	当 \(k\downarrow0\) 时，被积函数在闭区间上一致趋于 \(1\)，故 \(T\to2\pi\)。当 \(k\uparrow1\) 时，对每个 \(b<\pi/2\)，有
	\[
		\liminf_{k\uparrow1}T(h)
		\ge4\int_0^b\frac{\mathrm{d}\chi}{\cos\chi}.
	\]
	右侧在 \(b\uparrow\pi/2\) 时发散，所以周期趋于无穷。在 \(h=1\) 上，摆球到达上方平衡需要无穷时间；正文中 \(F^{-1}(1,0)\) 的非平衡轨道在两个时间方向都渐近该平衡，正对应振动周期的这一发散。
\end{solution}
